% \iffalse meta-comment
%%%%%%%%%%%%%%%%%%%%%%%%%%%%%%%%%%%%%%%%%%%%%%%%%
% fnttable.dtx
% Additions and changes Copyright (C) 2026 Clea F. Rees.
% Code from skeleton.dtx Copyright (C) 2015-2024 Scott Pakin (see below).
% Documentation and code from fonttable 
%   Copyright (C) 2005-2009 Peter R. Wilson
%                 2009-2011 Will Robertson
%                 2025-2026 LaTeX Project Team
% Documentation and code from unicodefonttable 
%   Copyright (C) 2019-2025 Frank Mittelbach
%
% This work may be distributed and/or modified under the
% conditions of the LaTeX Project Public License, either version 1.3c
% of this license or (at your option) any later version.
% The latest version of this license is in
%   https://www.latex-project.org/lppl.txt
% and version 1.3c or later is part of all distributions of LaTeX
% version 2008-05-04 or later.
%
% This work has the LPPL maintenance status 'muaintained'.
%
% The Current Maintainer of this work is Clea F. Rees.
%
% This work consists of all files listed in manifest.txt.
%
% The file fontscripts.dtx is a derived work under the terms of the
% LPPL. 
%
% The framework is based on version 2.4 of skeleton.dtx which is 
% part of dtxtut by Scott Pakin. A copy of dtxtut, including the
% unmodified version of skeleton.dtx is available from
% https://www.ctan.org/pkg/dtxtut and released under the LPPL.
%
% The content is based on version 2026-04-19 v1.6e of fonttable.dtx 
% which is part of fonttable by Peter R. Wilson (Herries Press) and
% 2025-07-11 v1.0k of unicodefonttable.dtx, which is part of 
% unicodefonttable by Frank Mittelbach. 
%
% A copy of fonttable, including the unmodified version of 
% fonttable.dtx is available from https://www.ctan.org/pkg/fonttable 
% and released under the LPPL.

% A copy of unicodefonttable, including the unmodified version of 
% unicodefonttable.dtx is available from 
% https://www.ctan.org/pkg/unicodefonttable and released under the LPPL.
%%%%%%%%%%%%%%%%%%%%%%%%%%%%%%%%%%%%%%%%%%%%%%%%%
%<sty|package>\NeedsTeXFormat{LaTeX2e}
\RequirePackage{svn-prov}
%<sty|package>\ProvidesPackageSVN[fnttable.sty]
%<driver>\ProvidesFileSVN
%<*driver|sty|package>
  {$Id: fnttable.dtx 12010 2026-08-27 01:46:27Z cfrees $}[v0.3 \revinfo][\filebase DTX: engine-agnostic fonttable utilities]
  \DefineFileInfoSVN[fnttable]
  \GetFileInfoSVN{fnttable}
%</driver|sty|package>
%
%<*driver>
\documentclass[twoside,10pt,british,]{ltxdoc}
\EnableCrossrefs
\CodelineIndex
\RecordChanges
%%^^A \OnlyDescription %% comment this out for the full glory
\usepackage{babel}
\usepackage{fancyhdr}
\pdfmapfile{-clm.map}
\pdfmapfile{+clm.map}
\usepackage[rm={proportional,lining},sf={proportional,lining},tt={monowidth,lining,tabular}]{cfr-lm}
\usepackage{array,tabularx}
\usepackage{xcolor}
\usepackage{xurl}
\urlstyle{sf}
\usepackage{microtype}
\usepackage{csquotes}
%%^^A addaswyd o chronos.tex
\MakeAutoQuote{‘}{’}
\MakeAutoQuote*{“}{”}
\usepackage{caption}
\DeclareCaptionFont{lf}{\lstyle}
\captionsetup[table]{labelfont=lf}
%%^^A sicrhau hyperindex=false: llwytho CYN bookmark
\usepackage{hypdoc}% ateb Ulrike Fischer: https://tex.stackexchange.com/a/695555/
\usepackage{bookmark}
\hypersetup{%
  colorlinks=true,
  citecolor={moss},
  extension=pdf,
  linkcolor={strawberry},
  linktocpage=true,
  pdfcreator={TeX},
  pdfproducer={pdfeTeX},
  urlcolor={blueberry}%
}
\usepackage{cleveref}
\ExplSyntaxOn
\NewDocumentCommand \ivals { +m }
{
  {
    \clist_if_empty:nF { #1 }
    {
      \normalfont
      \itshape
      < \clist_use:nn { #1 } { >\texttt{,} ~ < } >
    }
  }
}
\keys_define:nn { fnt / doc }
{
  unknown .code:n = {
    \cs_if_free:cT { \l_keys_key_str }
    {
      \tl_gset:cn { \l_keys_key_str } { #1 }
    }
  },
}
\NewDocumentCommand \fntdocset { +m }
{
  \keys_set:nn { fnt / doc } { #1 }
}
\ExplSyntaxOff
\fntdocset{%
  bug={\href{https://codeberg.org/cfr/nfssext/issues}{\textsc{bugtracker}}},
  codeberg={\href{https://codeberg.org/cfr/nfssext}{\textsc{codeberg}}},
  github={\href{https://github.com/cfr42/nfssext}{\textsc{github}}},
  ctan={\href{https://ctan.org/}{\textsc{ctan}}},
}
\newcommand*{\lpack}[1]{\textsf{#1}}
\title{\filebase}
\author{Clea F. Rees\thanks{%
    Bug tracker:
    \href{https://codeberg.org/cfr/nfssext/issues}{\url{codeberg.org/cfr/nfssext/issues}}
    \textbar{} Code:
    \href{https://codeberg.org/cfr/nfssext}{\url{codeberg.org/cfr/nfssext}}
    \textbar{} Mirror:
    \href{https://github.com/cfr42/nfssext}{\url{github.com/cfr42/nfssext}}%
}}
\date{\fileversion~\filedate}
\pagestyle{fancy}
\fancyhf{}
\fancyhf[lh]{\itshape\filebase}
\fancyhf[rh]{\itshape\fileversion}
\fancyhf[cf]{\itshape--- \thepage~/~\lastpage{} ---}
\ExplSyntaxOn
\hook_gput_code:nnn {shipout/lastpage} {.}
{
  \property_record:nn {t:lastpage}{abspage,page,pagenum}
}
\cs_new_protected_nopar:Npn \lastpage 
{
  \property_ref:nn {t:lastpage}{page}
}
\NewDocumentCommand \plarg {+m} {{\ttfamily (\ivals{#1})}}
\ExplSyntaxOff
\definecolor{strawberry}{rgb}{1.000,0.000,0.502}
\definecolor{blueberry}{rgb}{0.000,0.000,1.000}
\definecolor{moss}{rgb}{0.000,0.502,0.251}
\makeatletter
  \def\@xobeysp{\leavevmode\penalty100\ }
\makeatother
%%^^A\setcounter{StandardModuleDepth}{1}
%%^^A\makeatletter
%%^^A  \@mparswitchfalse
%%^^A\makeatother
%%^^A\newenvironment{addtomargins}[1]{%
%%^^A  \begin{list}{}{%
%%^^A    \topsep 0pt%
%%^^A    \addtolength{\leftmargin}{#1}%
%%^^A    \addtolength{\rightmargin}{#1}%
%%^^A    \listparindent \parindent
%%^^A    \itemindent \parindent
%%^^A    \parsep \parskip}%
%%^^A  \item[]}{\end{list}}
\begin{document}
  %%^^A \raggedbottom
  \DocInput{\filename}
\end{document}
%</driver>
%
% \fi
%
%
% \DoNotIndex{\',\.,\@M,\@@input,\@addtoreset,\@arabic,\@badmath}
% \DoNotIndex{\@centercr,\@cite}
% \DoNotIndex{\@dotsep,\@empty,\@float,\@gobble,\@gobbletwo,\@ignoretrue}
% \DoNotIndex{\@input,\@ixpt,\@m}
% \DoNotIndex{\@minus,\@mkboth,\@ne,\@nil,\@nomath,\@plus,\@set@topoint}
% \DoNotIndex{\@tempboxa,\@tempcnta,\@tempdima,\@tempdimb}
% \DoNotIndex{\@tempswafalse,\@tempswatrue,\@viipt,\@viiipt,\@vipt}
% \DoNotIndex{\@vpt,\@warning,\@xiipt,\@xipt,\@xivpt,\@xpt,\@xviipt}
% \DoNotIndex{\@xxpt,\@xxvpt,\\,\ ,\addpenalty,\addtolength,\addvspace}
% \DoNotIndex{\advance,\Alph,\alph}
% \DoNotIndex{\arabic,\ast,\begin,\begingroup,\bfseries,\bgroup,\box}
% \DoNotIndex{\bullet}
% \DoNotIndex{\cdot,\cite,\CodelineIndex,\cr,\day,\DeclareOption}
% \DoNotIndex{\def,\DisableCrossrefs,\divide,\DocInput,\documentclass}
% \DoNotIndex{\DoNotIndex,\egroup,\ifdim,\else,\fi,\em,\endtrivlist}
% \DoNotIndex{\EnableCrossrefs,\end,\end@dblfloat,\end@float,\endgroup}
% \DoNotIndex{\endlist,\everycr,\everypar,\ExecuteOptions,\expandafter}
% \DoNotIndex{\fbox}
% \DoNotIndex{\filedate,\filename,\fileversion,\fontsize,\framebox,\gdef}
% \DoNotIndex{\global,\halign,\hangindent,\hbox,\hfil,\hfill,\hrule}
% \DoNotIndex{\hsize,\hskip,\hspace,\hss,\if@tempswa,\ifcase,\or,\fi,\fi}
% \DoNotIndex{\ifhmode,\ifvmode,\ifnum,\iftrue,\ifx,\fi,\fi,\fi,\fi,\fi}
% \DoNotIndex{\input}
% \DoNotIndex{\jobname,\kern,\leavevmode,\let,\leftmark}
% \DoNotIndex{\list,\llap,\long,\m@ne,\m@th,\mark,\markboth,\markright}
% \DoNotIndex{\month,\newcommand,\newcounter,\newenvironment}
% \DoNotIndex{\NeedsTeXFormat,\newdimen}
% \DoNotIndex{\newlength,\newpage,\nobreak,\noindent,\null,\number}
% \DoNotIndex{\numberline,\OldMakeindex,\OnlyDescription,\p@}
% \DoNotIndex{\pagestyle,\par,\paragraph,\paragraphmark,\parfillskip}
% \DoNotIndex{\penalty,\PrintChanges,\PrintIndex,\ProcessOptions}
% \DoNotIndex{\protect,\ProvidesClass,\raggedbottom,\raggedright}
% \DoNotIndex{\refstepcounter,\relax,\renewcommand,\reset@font}
% \DoNotIndex{\rightmargin,\rightmark,\rightskip,\rlap,\rmfamily,\roman}
% \DoNotIndex{\roman,\secdef,\selectfont,\setbox,\setcounter,\setlength}
% \DoNotIndex{\settowidth,\sfcode,\skip,\sloppy,\slshape,\space}
% \DoNotIndex{\symbol,\the,\trivlist,\typeout,\tw@,\undefined,\uppercase}
% \DoNotIndex{\usecounter,\usefont,\usepackage,\vfil,\vfill,\viiipt}
% \DoNotIndex{\viipt,\vipt,\vskip,\vspace}
% \DoNotIndex{\wd,\xiipt,\year,\z@}
% \DoNotIndex{\verb,\ProvidesPackageSVN,\NeedsTeXFormat,\ProcessKeyOptions}
%
%<*skip>
% \changes{v1.0}{2005/11/27}{First public release}
% \changes{v1.0a}{2005/12/06}{Minor bug fix}
% \changes{v1.1}{2006/10/02}{Minor extensions}
% \changes{v1.2}{2006/10/02}{Minor extensions of texts}
% \changes{v1.3}{2009/04/30}{Fix serious bug}
% \changes{v1.4}{2009/05/06}{Added another way to select a font}
% \changes{v1.5}{2009/05/12}{Added code for Knuth's testfont.tex}
% \changes{v1.51}{2009/05/14}{Eliminated a clash with babel package}
% \changes{v1.5b}{2009/09/02}{New maintainer (Will Robertson)}
% \changes{v1.6}{2009/10/15}{New spacing of the decimals from Peter Wilson}
% \changes{v1.6b}{2011/02/13}{Fix bug with Spanish babel}
% \changes{v1.6c}{2017/06/06}{Fix bug with fonts with `at' in their name}
% \changes{v1.6d}{2025/03/12}{Missing relax (issue 51)}
% \changes{v1.6e}{2026/04/19}{Avoid global definitions (https://tex.stackexchange.com/questions/762029)}
%</skip>
%
% \newcommand*{\Lopt}[1]{\textsf {#1}}            ^^A typeset an option
% \newcommand*{\file}[1]{\texttt {#1}}            ^^A typeset a file
% \newcommand*{\Lcount}[1]{\textsl {\small#1}}    ^^A typeset a counter
% \newcommand*{\pstyle}[1]{\textsl {#1}}          ^^A typeset a pagestyle
% \newcommand*{\Lenv}[1]{\texttt {#1}}            ^^A typeset an environment
% \newcommand*{\AD}{\textsc{ad}}
% \newcommand*{\thisfont}{OandS}
%
% \maketitle
% \begin{abstract}
%   The package lets you typeset the characters in a font in tabular and/or
%     running text forms.
%   It aims to provide features at least equivalent to those provided by
%     \lpack{fonttable}, for both 8bit and Unicode engines\footnote{%
%       Where Unicode support is not required, \lpack{fonttable} should be
%         preferred, since it is maintained by the \LaTeX{} Project.
%       For Unicode engines, however, there is only \lpack{unicodefonttable}, 
%         which is suitable \emph{only} for tables of characters.
%       Moreover, while \lpack{fonttable} works only with non-Unicode engines,
%         \lpack{unicodefonttable} works only with Unicode ones.
%       But the biggest problem, for my purposes, is that it loads 
%         \lpack{fontspec} unconditionally \dots
%     }.
%   At present, however, only running forms are supported for Unicode\footnote{%
%     \lpack{unicodefonttable} may be used to produce font tables for these
%       engines.
%       Unfortunately, the interface and features differ from that provided by 
%         \lpack{fonttable}, which is the \emph{raison d'être} for this package.%
%   }.
%
%   \textbf{Note that this package assumes \emph{only} format-level support
%     for font configuration i.e.~NFSS plus \lpack{luaotfload} on Unicode 
%     engines.
%     It does \emph{not} load or support \lpack{fontspec}.%
%   }
% \end{abstract}
% \tableofcontents
%
% \section{Introduction}
%
%    The \lpack{fnttable} package lets you typeset a font's character set
% in tabular and/or running text forms.
%
% \lpack{fnttable} is intended as a \emph{replacement} for 
% \lpack{fonttable}, when support for Unicode fonts is required.
% \textbf{One cannot load both \lpack{fnttable} \emph{and}
% \lpack{fonttable} in the same document!}
%
% This manual is typeset according to the conventions of the
% \LaTeX{} \textsc{docstrip} utility which enables the automatic
% extraction of the \LaTeX{} macro source files~\cite{COMPANION}.
%
% \section{The package}
%
%    The package provides commands to typeset a table of all the glyphs in
% a given font and to typeset an example of regular text. For font designers
% it provides commands to typeset a `test' glyph among sets of glyphs
% from the font.
%
% \DescribeMacro{\fnthours}
% As a convenience, \cs{fnthours} prints the time of day when the file was 
% processed; it uses the 24 hour clock notation. (The macro \cs{today} prints
% the date when the file was processed.)
%
% \subsection{Table and texts}
%
%
% \DescribeMacro{\fonttable}
% The command \\
% \cs{fonttable}\marg{testfont} \\
% typesets a table showing all the
% glyphs in the \textbf{7bit} or \textbf{8bit} \meta{testfont}, where 
% \meta{testfont} is the name of a 
% font file\footnote{More precisely, the name of a \texttt{.tfm} file.}
% like \texttt{cmr10} (for Computer Modern Roman) or \texttt{pzdr} 
% (for Zapf Dingbats).
%
% NOTE: The \lpack{mftinc} package~\cite{MFTINC} for pretty-printing 
% METAFONT code also defines a \cs{fonttable} macro that is akin to this 
% one. If you want to use both packages together then you can use the following
% general procedure for when a macro \cs{macro} is defined in both 
% \lpack{packA} and \lpack{packB} packages.
% \begin{verbatim}
% \usepackage{packA}
% \let\macroA\macro%   save packA's definition
% \let\macro\relax%    undefine \macro
% \usepackage{packB}%  now it's packB's definition of \macro
% ...
% \macro %   use the packB defintion
% \macroA %  use the packA definition   
% \end{verbatim}
%
% \DescribeMacro{\xfonttable}
% The command \\
% \cs{xfonttable}\marg{encoding}\marg{family}\marg{series}\marg{shape} \\
% typesets a table showing all the
% glyphs in the font with \textbf{7} or \textbf{8bit} encoding 
% \meta{encoding} (e.g., T1 or OMS),
% family \meta{family} (e.g., \texttt{ppl} for Palatino or \texttt{cmbrs}
% for CM Bright Math (OMS)), font series \meta{series} (e.g., \texttt{sb}
% for semibold of \texttt{m} for medium), and font shape \meta{shape}
% (e.g., \texttt{n} for normal or \texttt{sc} for small caps).
% For example: \\
% \verb?\xfonttable{U}{pzd}{m}{n}? \\
% for Zapf Dingbats.
%
% \DescribeMacro{\pikfont}
% The command\footnote{The name was chosen in an attempt to avoid
% clashes with other macros that might perform similar functions.}   \\
% \cs{pikfont}\marg{encoding}\marg{family}\marg{series}\marg{shape} \\
% selects the font with encoding \meta{encoding} (e.g., T1, TU or OMS),
% family \meta{family} (e.g., \texttt{ppl} for Palatino or \texttt{cmbrs}
% for CM Bright Math (OMS)), font series \meta{series} (e.g., \texttt{sb}
% for semibold of \texttt{m} for medium), and font shape \meta{shape}
% (e.g., \texttt{n} for normal or \texttt{sc} for small caps).
% For example: \\
% \verb?\pikfont{T1}{ppl}{m}{sc}? \\
% for Palatino small caps.
% The size of the font corresponds to the current setting (e.g.,
% \cs{footnotesize}, \cs{normalsize}, \cs{Large}). It can also be
% changed after being selected by the incantation \\
% \cs{fontsize}\marg{size}\marg{baselineskip}\cs{selectfont} \\
% where \meta{size} is the normal height and \meta{baselineskip} is the
% distance between text lines; the measurement system is \texttt{pts}
% but just use numbers with no units specified. For example: \\
% \verb?\fontsize{12}{15}\selectfont? \\
% for a 12pt font with 15pts between baselines.
%
% If you are unsure about the meaning of the various arguments 
% of \cs{xfonttable} and \cs{pikfont} see
% \textit{The Companion}~\cite[Chapter 7]{COMPANION} or the 
% \textit{LaTeX2e font selection} manual (\texttt{fntguide.tex}; try
% \verb?texdoc fntguide?).
%
% \DescribeMacro{\fontrange}
% The package attempts to populate the table with a maximum of 256 glyphs,
% numbered from 0 to 255.
% The \cs{fontrange}\marg{low}\marg{high} declaration changes this by reducing
% the range so that it extends from \meta{low} to \meta{high}, where \meta{low}
% should be at least 0 and \meta{high} at most 256, and \meta{low} less than
% \meta{high}.
%
%    The table is composed of blocks of sixteen characters. If necessary the
% value of \meta{low} is adjusted lower and \meta{high} is adjusted higher
% to match this block structure. For example, if you wanted a table of the
% lower 128 characters then \verb?\fontrange{0}{127}? would do the job,
% while the upper half of a 256 character font could be tabulated via
% \verb?\fontrange{128}{255}?.
%
% \DescribeMacro{\decimals}
% \DescribeMacro{\nodecimals}
% Normally each cell in the table includes the decimal number of the position
% in the (256) character set. \cs{nodecimals} turns off this numbering and
% \cs{decimals} turns it on. The default is \cs{decimals}.
%
% \DescribeMacro{\hexoct}
% \DescribeMacro{\nohexoct}
% Normally the columns and rows in the table are numbered using hexadecimal 
% and octal numbers. These can be turned off by \cs{nohexoct} and turned 
% on again with \cs{hexoct}, which is the default.
%
% \DescribeMacro{\ftablewidth}
% The font table's width is the length \cs{ftablewidth}, which by default
% is set to the normal textwidth (or more exactly, to \cs{hsize}). The table
% itself is left aligned. However, if \cs{nohexoct} is in effect the width of
% the table is its natural width.
%
% \DescribeMacro{\fntcolwidth}
% When \cs{nohexoct} is in effect the minimum width of a table column
% is \cs{fntcolwidth}. This is initially declared as \\
% \verb?\setwidth{\fntcolwidth}{0.08\ftablewidth}?
%
% \DescribeMacro{\fonttext}
% The command \cs{fonttext}\marg{testfont} typesets an example text using
% the \meta{testfont} (e.g. \texttt{cmr10}).
%
% \DescribeMacro{\simpletext}
% \DescribeMacro{\fulltext}
% The example text can be just a paragraph and a line of capitals,
%  or include more complex accented
% words as well. Following the declaration \cs{fulltext} the complex words are
% included as well as the example paragraph. The default is \cs{simpletext}
% for just the paragraph.
%
% \DescribeMacro{\regulartext}
% The command \cs{regulartext}\marg{fontspec} typesets the example text
% using \meta{fontspec}, for example \verb?\rmfamily\itshape?
% or \verb?\pikfont{T1}{pnc}{m}{it}?.
%
% \DescribeMacro{\fonttexts}
% \DescribeMacro{\regulartexts}
% The macro \cs{fonttexts}\marg{testfont}\marg{text} typesets \meta{text} using
% the \meta{testfont} (e.g., \texttt{cmr10}). Similarly the macro
% \cs{regulartexts}\marg{fontspec}\marg{text} typesets \meta{text}
% using \meta{fontspec} (e.g., \verb?\rmfamily\itshape? or 
% \verb?\pikfont{T1}{ppl}{m}{it}?).
%
% \DescribeMacro{\germanparatext}
% \DescribeMacro{\latinparatext}
% \cs{germanparatext} expands to a German language paragraph, borrowed from the 
% \lpack{blindtext} package~\cite{BLINDTEXT}. 
% \cs{latinparatext} expands to one version of
% a paragraph of the traditional \textit{lorem ipsum} dummy Latin text.
% Either, or both, of these could be used as the \meta{text} argument
% to \cs{fonttexts} or \cs{regulartexts}. \\
% NOTE: These were originally called \cs{germantext} and \cs{latintext}
%  but on 2009/05/14 I was told that the \lpack{babel} package defines \cs{latintext}, 
%  which causes unexpected results if it is used in the same document as this
%  package. To try and be on the safe side I renamed \cs{germantext}
% as well as \cs{latintext}.
%
% \DescribeMacro{\aztext}
% \DescribeMacro{\AZtext}
% \DescribeMacro{\digitstext}
% \DescribeMacro{\punctext}
% \cs{aztext} expands to the lowercase Latin alphabet a to z, and \cs{AZtext}
% is the corresponding command for the uppercase A to Z. 
% The macros \cs{digitstext} and \cs{punctext} expand respectively to the
% digits 0 to 9, and to the typical punctuation marks. In all cases there 
% is a space between each character. 
%
% \subsection{Testing a glyph}
%
%    The macros here are a reimplementation of Donald Knuth's 
% \texttt{testfont.tex}, which is available from CTAN. 
%
%    In the following, the value of a glyph argument can be specified as
% its location in the font (i.e., as a decimal number). With a few exceptions,
% if the glyph is within the visible ASCII range (33--126) it may instead be 
% specified by the ASCII character prefixed with a single open quote 
% mark\footnote{Sometimes called a
% `backquote'.} (`). The exceptions are nos: 37 (\verb?%?), 92 (\verb?\?)
% 123 (\verb?{?) and 125 (\verb?}?) (but there may be others). 
% In any case, the glyph representing the
% character \texttt{p} can be specified either as \texttt{`p} or as 
% \texttt{112}. 
%
%    The glyphs are taken from the current font. If the font does not
% have Latin alphabet glyphs in the ASCII locations then 
% in the descriptions below phrases like `lowercase alphabet' or `uppercase
% alphabet' or `digits', should be taken to mean (the glyphs in) those 
% locations.
%
% \DescribeMacro{\glyphmixture}
% \cs{glyphmixture}\marg{T}\marg{S}\marg{E} typesets the \meta{T} (test) glyph 
% between the glyphs in the range from \meta{S} (start) to \meta{E} (end).
% For example \\
% \verb?\glyphmixture{`e}{`f}{`g}? will produce \\
% \texttt{efeeffeeefffef} \\
% \texttt{egeeggeeegggeg}
%
% \DescribeMacro{\glyphalternation}
% \cs{glyphalternation}\marg{T}\marg{S}\marg{E} typesets the \meta{T} 
% glyph alternately between each glyph in the range from \meta{S} 
% to \meta{E}.
% For example \\
% \verb?\glyphalternation{`e}{`f}{`g}? will produce \\
% \texttt{efefefefefefefefe} \\
% \texttt{egegegegegegegege}
%
% \DescribeMacro{\glyphseries}
% \cs{glyphseries}\marg{T}\marg{S}\marg{E} typesets the \meta{T}
% glyph between the glyphs in the range from \meta{S}
% to \meta{E}.
% For example \\
% \verb?\glyphseries{`e}{`f}{`h}? will produce \\
% \texttt{efegehe}
%
% \DescribeMacro{\glyphalphabet}
% \DescribeMacro{\GLYPHALPHABET}
% \cs{glyphalphabet}\marg{T} typesets the \meta{T} glyph between
% each letter of the lowercase Latin alphabet plus a few others. 
% \cs{GLYPHALPHABET}\marg{T}
% does the same but using the uppercase Latin alphabet.
% For example, the output of \\
% \verb?\glyphalphabet}{`3}? is like \\
% \texttt{3a3b3c3d3e3f3g\ldots3z3\char31 3}\verb?~?\texttt{3\char33 3\char34 3}
%
% \DescribeMacro{\glyphlowers}
% \DescribeMacro{\glyphlowers}
% \DescribeMacro{\glyphdigits}
% \cs{glyphlowers} takes each character of the lowercase alphabet in turn as
% a test glyph and sets it interspersed among the other lowercase characters.
% \cs{glyphuppers} and \cs{glyphdigits} are similar except that they use
% the uppercase alphabet and the ten digits instead. For example, 
% \cs{glyphdigits} produces output like \\
% \texttt{000102030405060708090} \\
% \texttt{101112131415161718191} \\
% \texttt{202122232425262728292} \\
% \ldots \\
% \texttt{909192939495969798999}
%
% \DescribeMacro{\glyphpunct}
% \cs{glyphpunct} sets a collection of words with an assortment of
% punctuation marks.
%
%
% \StopEventually{
% \bibliographystyle{alpha}
% \renewcommand{\refname}{Bibliography}
% \begin{thebibliography}{GMS94}
% \addcontentsline{toc}{section}{\refname}
%
% \bibitem[Lik05]{BLINDTEXT}
% Knut Lickert.
% \newblock \emph{Blindtext.sty: Creating text for testing / Texterzeugung
%                 zum testen}.
% \newblock October 2005.
% \newblock (Available from CTAN in \url{macros/latex/contrib/blindtext})
%
% \bibitem[MG04]{COMPANION}
% Frank Mittelbach and Michel Goossens.
% \newblock \emph{The LaTeX Companion}.
% \newblock Second edition.
% \newblock Addison-Wesley Publishing Company, 2004.
%
% \bibitem[Pak05]{MFTINC}
% Scott Pakin
% \newblock `The \lpack{mftinc} package',
% \newblock January 2005.
% \newblock (Available from CTAN in \url{macros/latex/contrib/mftinc})
%
%
% \end{thebibliography}
%
% \PrintIndex
%
% }
%
%
%
% \section{The code} \label{sec:mf}
%
%    \begin{macrocode}
%<*sty>
%<@@=fnttable>
\ExplSyntaxOn
\msg_new:nnnn {fnttable}{incomp} {
  You~tried~to~load~`fnttable'~and~`#1'~\msg_line_context:,~
  \iow_newline: but~these~are~not~compatible.
}{`fnttable',~`fonttable'~and~`unicodefonttable'~are~incompatible~
  \io_newline: alternatives.~Load~at~most~one!
}
\AddToHook{pkg/fonttable/before}{
  \msg_error:nnn {fnttable}{incomp}{fonttable}
}
\AddToHook{pkg/unicodefonttable/before}{
  \msg_error:nnn {fnttable}{incomp}{unicodefonttable}
}
\ExplSyntaxOff
%    \end{macrocode}
%
% \subsection{Table and texts}
%
% Most of the code below is an edited version of code used in 
% \texttt{nfssfont.tex} for displaying aspects of the set of glyphs in a font.
%
% \begin{macro}{\sevenrm}
% A small fixed size roman font.
%    \begin{macrocode}
\providecommand*{\sevenrm}{\fontsize{7}{9pt}\rmfamily}
%    \end{macrocode}
% \end{macro}
%
% \begin{macro}{\f@tm}
% \begin{macro}{\f@tn}
% \begin{macro}{\f@tp}
% \begin{macro}{\f@tmax}
% \begin{macro}{\f@tdim}
% \begin{macro}{\f@tuni}
% Counts, a dimen and a bool.
%    \begin{macrocode}
\newcount\f@tm \newcount\f@tn \newcount\f@tp \newdimen\f@tdim
\newcount\f@tmax
\f@tmax\@cclvi
\newif\iff@tuni

%    \end{macrocode}
% \end{macro}
% \end{macro}
% \end{macro}
% \end{macro}
% \end{macro}
%
% \begin{macro}{\f@t@uni@check}
% If Unicode encoding do one thing; if not, do another.
%    \begin{macrocode}
\newcommand\iff@t@uni@check[3]{%
  \edef\tempa{#1}\edef\tempb{TU}\edef\tempc{EU1}\edef\tempf{EU2}%
  \f@tunifalse
  \ifx\tempa\tempb
    \f@tunitrue
  \else
    \ifx\tempa\tembc
      \f@unitrue
    \else
      \ifx\tempa\tempd
        \f@unitrue
      \fi
    \fi
  \fi
  \iff@tuni
    #2
  \else
    #3
  \fi
}
%
%    \end{macrocode}
% \end{macro}
%
% \begin{macro}{\f@t@uni@error}
% If Unicode encoding do one thing; if not, do another.
%    \begin{macrocode}
\newcommand\f@t@uni@error{%
  \PackageError{fnttable}{Sorry, tables are not yet supported for
    Unicode encodings.%
  }{%
    Consider `unicodefonttable` if you require tables for Unicode fonts.%
  }%
}
%
%    \end{macrocode}
% \end{macro}
%
% \begin{macro}{\fonttable}
% \cs{fonttable}\marg{font} typesets a table of all the glyphs in the 
% \meta{font} (e.g., auncl10).
%    \begin{macrocode}
\newcommand*{\fonttable}[1]{%
  \iff@t@uni@check{TU}{\f@t@uni@error}{%
    \def\f@tfontname{#1}%
    \bgroup
    \f@tstartfont 
    \ftable
    \egroup
  }%
}

%    \end{macrocode}
% \end{macro}
%
% \begin{macro}{\pikfont}
% \cs{pikfont}\marg{encoding}\marg{family}\marg{series}\marg{shape} 
% selects the font with \meta{encoding},
% \meta{family}, \meta{series} and \meta{shape}.
% \changes{v1.4}{2009/05/06}{Added \cs{pikfont}}
%    \begin{macrocode}
\DeclareRobustCommand{\pikfont}[4]{%
  \fontencoding{#1}\fontfamily{#2}\fontseries{#3}\fontshape{#4}\selectfont
}

%    \end{macrocode}
% \end{macro}
%
% \begin{macro}{\xfonttable}
% \cs{xfonttable}\marg{encoding}\marg{family}\marg{series}\marg{shape} 
% typesets a table of all the glyphs in the font with \meta{encoding},
% \meta{family}, \meta{series} and \meta{shape}
% (e.g., \verb?\xfonttable{T1}{pnc}{m}{it}? for New Century Schoolbook italic).
% The original code for the macro was supplied by Enrico Gregorio.
% \changes{v1.4}{2009/05/06}{Added \cs{xfonttable}}
%    \begin{macrocode}
\newcommand*{\xfonttable}[4]{%
  \begingroup
    \pikfont{#1}{#2}{#3}{#4}%
    \edef\f@tfontname{\fontname\font}%
%    \end{macrocode}
% New: strip any size information from the fontname (which could be, e.g.,
% either `\verb*|cmr10|' or `\verb*|cmr10 at 10pt|'.)
% This wasn't necessary before because we didn't explicitly choose the font
% size; it was inferred automatically.
%    \begin{macrocode}
    \edef\@tempa{ \string a\string t}%
    \edef\@tempb{\noexpand\in@{\@tempa}{\f@tfontname}}%
    \@tempb
    \ifin@
      \edef\f@tfontname{\expandafter\f@tstripsize\f@tfontname}%
    \fi
%    \end{macrocode}
% End new code, and finish as before:
%    \begin{macrocode}
    \normalfont
    \f@tstartfont
    \iff@t@uni@check{#1}{%
      \f@tmax=10000
      \let\f@ttestrow\f@t@uni@testrow
      \f@t@uni@ftable
    }{%
      \f@tmax\@cclvi
      \ftable
    }%
  \endgroup
}
%    \end{macrocode}
% \end{macro}
%
%
% \begin{macro}{\f@tstripsize}
% Needed above.
%    \begin{macrocode}
\edef\@tempa{%
  \def\noexpand\f@tstripsize
    ##1 \string a\string t##2\string p\string t{##1}%
}
\@tempa
%    \end{macrocode}
% \end{macro}
%
%
% \begin{macro}{\f@tstartfont}
% Sets up for a font table.
%    \begin{macrocode}
\newcommand*{\f@tstartfont}{%
%    \end{macrocode}
% New: scale the font by 0.01\% to (attempt to) avoid TeX's font optimisation.
% This becomes a problem in Spanish babel, say, when \verb|\textfont\fam|
% changes when \verb|cmr10| has been loaded under a different name, here.
% (And the \verb|\textfont| can no longer be parsed correctly. See:
% \url{http://latex-alive.tumblr.com/post/3229118083/texs-font-loading-optimisation}
% )
%    \begin{macrocode}
  \@tempdima=\f@size pt
  \font\f@ttestfont=\f@tfontname\space at 0.9999\@tempdima\relax
%    \end{macrocode}
% Continue as before:
%    \begin{macrocode}
  \f@ttestfont \f@tsetbaselineskip
  \ifdim\fontdimen6\f@ttestfont<10pt\relax
    \rightskip=0pt plus 20pt\relax
  \else
    \rightskip=0pt plus 2em\relax
  \fi
  \spaceskip=\fontdimen2\f@ttestfont % space between words (\raggedright)
  \xspaceskip=\fontdimen2\f@ttestfont
  \advance\xspaceskip by\fontdimen7\f@ttestfont
}
%    \end{macrocode}
% \end{macro}
%
%
%
% \begin{macro}{\f@tsetbaselineskip}
%    \begin{macrocode}
\newcommand*{\f@tsetbaselineskip}{%
  \setbox0=\hbox{%
    \f@tn=0
    \loop\char\f@tn \ifnum \f@tn<255 \advance\f@tn 1 \repeat
  }
  \baselineskip=6pt 
  \advance\baselineskip\ht0 
  \advance\baselineskip\dp0 
}

%    \end{macrocode}
% \end{macro}
%
% \begin{macro}{\f@toct}
% \cs{f@toct}\marg{onum} typesets the octal constant \meta{onum}.
%    \begin{macrocode}
\newcommand*{\f@toct}[1]{\hbox{\rmfamily\'{}\kern-.2em\itshape
           #1\/\kern.05em}} % octal constant
%    \end{macrocode}
% \end{macro}
%
% \begin{macro}{\f@thex}
% \cs{f@thex}\marg{hnum} typesets the hexadecimal constant \meta{hnum}.
%    \begin{macrocode}
\newcommand*{\f@thex}[1]{\hbox{\rmfamily\H{}\ttfamily#1}} % hexadecimal constant
%    \end{macrocode}
% \end{macro}
%
% \begin{macro}{\f@tsetdigs}
% \cs{f@tsetdigs}
%    \begin{macrocode}
\def\f@tsetdigs#1"#2{\gdef\f@thexpre{#2}% \f@thexpre=hex prefix; \f@tocta\f@toctb=corresponding octal
 \f@tm=\f@tn \divide\f@tm by 64 \xdef\f@tocta{\the\f@tm}%
 \multiply\f@tm by-64 \advance\f@tm by\f@tn \divide\f@tm by 8 \xdef\f@toctb{\the\f@tm}}
%    \end{macrocode}
% \end{macro}
%
% \begin{macro}{\f@ttestrow}
% \cs{f@ttestrow} checks if there are any characters in the next block of 
% 16 slots.
%    \begin{macrocode}
 \newcommand*{\f@ttestrow}{\setbox0=\hbox{\penalty 1\def\\{\char"\f@thexpre}%
 \\0\\1\\2\\3\\4\\5\\6\\7\\8\\9\\A\\B\\C\\D\\E\\F%
 \global\f@tp=\lastpenalty}} % \f@tp=1 if none of the characters exist

%    \end{macrocode}
% \end{macro}
%
% \begin{macro}{\f@t@uni@testrow}
% \cs{f@ttestrow} checks if there are any characters in the next block of 
% 16 slots.
% TODO!
%    \begin{macrocode}
 \newcommand*{\f@t@uni@testrow}{%
   \setbox0=\hbox{%
     \penalty 1%
     \def\\{\char"\f@thexpre}%
     \\0\\1\\2\\3\\4\\5\\6\\7\\8\\9\\A\\B\\C\\D\\E\\F%
     \global\f@tp=\lastpenalty
   }%
} % \f@tp=1 if none of the characters exist

%    \end{macrocode}
% \end{macro}
%
% \begin{macro}{\ifhexoct}
% \begin{macro}{\hexoct}
% \begin{macro}{\nohexoct}
% Flag for (not) setting hex and octal numbers.
%    \begin{macrocode}
\newif\ifhexoct
\newcommand*{\hexoct}{\hexocttrue}
\newcommand*{\nohexoct}{\hexoctfalse}
\hexoct

%    \end{macrocode}
% \end{macro}
% \end{macro}
% \end{macro}
% 
%
% \begin{macro}{\f@toddlinenum}
% \cs{f@toddline}
%    \begin{macrocode}
\newcommand*{\f@toddline}{\cr
  \noalign{\nointerlineskip}
  \multispan{19}\hrulefill&
  \setbox0=\hbox{\lower 2.3pt\hbox{\f@thex{\f@thexpre x}}}\smash{\box0}
  \cr
  \noalign{\nointerlineskip}}

%    \end{macrocode}
% \end{macro}
%
% \begin{macro}{\iff@tskipping}
% \begin{macro}{\f@tskippingtrue}
% \begin{macro}{\f@tskippingfalse}
%    \begin{macrocode}
\newif\iff@tskipping

%    \end{macrocode}
% \end{macro}
% \end{macro}
% \end{macro}
%
% \begin{macro}{\fontrange}
% \cs{fontrange}\marg{low}\marg{high} sets the character range to be 
% output.
%    \begin{macrocode}
\newcommand*{\fontrange}[2]{%
  \ifnum#1<#2\relax
%    \end{macrocode}
% Set \cs{f@tlow} to the nearest multiple of 16 that is at or below \meta{low},
% but first make sure that it will be at least 0.
%    \begin{macrocode}
  \ifnum#1<\z@
    \f@tm=\z@
  \else
    \f@tm=#1
    \divide \f@tm \sixt@@@@n
    \multiply \f@tm \sixt@@@@n
  \fi
  \edef\f@tlow{\the\f@tm}
%    \end{macrocode}
% Set \cs{f@thigh} to the nearest multiple of 16 at or above \meta{high},
% finally making sure that its maximum is 256.
%    \begin{macrocode}
  \f@tm=#2
  \divide \f@tm \sixt@@@@n
  \advance \f@tm \@ne
  \multiply \f@tm \sixt@@@@n
  \ifnum \f@tm > \f@tmax  \f@tm \f@tmax \fi
  \edef\f@thigh{\the\f@tm}
 \else
   \PackageError{fnttable}{%
      Improper values for fontrange. Default values substituted}{\@ehc}
   \def\f@tlow{0} \def\f@thigh{256}
  \fi
}
\fontrange{0}{256}

%    \end{macrocode}
% \end{macro}
%
% \begin{macro}{\f@tloopforsixteen}
% \cs{f@tloopforsixteen} sets up a block of sixteen character slots.
%    \begin{macrocode}
\newcommand*{\f@tloopforsixteen}{%
  \ifnum\f@tn<\f@tlow \global\f@tn=\f@tlow\fi
  \loop\f@tskippingfalse
  \ifnum\f@tn<\f@thigh \f@tm=\f@tn \divide\f@tm \sixt@@@@n \chardef\next=\f@tm
  \expandafter\f@tsetdigs\meaning\next \f@ttestrow
  \ifnum\f@tp=\@ne \f@tskippingtrue \fi\fi
  \iff@tskipping \global\advance\f@tn \sixt@@@@n \repeat}

%    \end{macrocode}
% \end{macro}
%
% \begin{macro}{\f@tevenline}
% \begin{macro}{\f@tevenlinenonum}
% \cs{f@tevenline} gets next non-empty set of a block of 16 characters.
% It either calls \cs{f@tmorechart} to print them, or \cs{f@tendchart} to
% finish off the table if all 256 potential characters have been processed.
%
% \cs{f@tevenlinenonum} does something similar when no external numbers
% are printed.
%    \begin{macrocode}
\newcommand*{\f@tevenline}{%
  \f@tloopforsixteen
  \ifnum\f@tn=\f@thigh \let\next=\f@tendchart\else\let\next=\f@tmorechart\fi
  \next}
\newcommand*{\f@tevenlinenonum}{%
  \f@tloopforsixteen
  \ifnum\f@tn=\f@thigh\relax
    \\\hline
  \else
    \\\hline
    \f@tmorechartnonum
  \fi}

%    \end{macrocode}
% \end{macro}
% \end{macro}
%
% \begin{macro}{\f@tmorechart}
% \begin{macro}{\f@tmorechartnonum}
% \cs{f@tmorechart} sets two lines of the table, and \cs{f@tmorechartnonum}
% does the same when there are no external numbers.
%    \begin{macrocode}
\newcommand*{\f@tmorechart}{\cr\noalign{\hrule\penalty5000}
 \f@tchartline \f@toddline \f@tm=\f@toctb \advance\f@tm 1 \xdef\f@toctb{\the\f@tm}
 \f@tchartline \f@tevenline}
\newcommand*{\f@tmorechartnonum}{%
  \f@tsimpleline \\ \hline
  \f@tsimpleline \f@tevenlinenonum}

%    \end{macrocode}
% \end{macro}
% \end{macro}
%
% \begin{macro}{\f@tchartline}
% \begin{macro}{\f@tsimpleline}
% \cs{f@tchartline} does a line of the table, including external numbers,
%  and \cs{f@tsimpleline} does an unnumbered line.
%    \begin{macrocode}
\newcommand*{\f@tchartline}{%
  &\f@toct{\f@tocta\f@toctb x}&&\f@tpsg{}&&\f@tpsg{}&&\f@tpsg{}&&\f@tpsg{}&&\f@tpsg{}&&\f@tpsg{}&&\f@tpsg{}&&\f@tpsg{}&&}
\newcommand*{\f@tsimpleline}{%
  \f@tpsg{}\f@tchartstrut& \f@tpsg{} & \f@tpsg{} & \f@tpsg{} & \f@tpsg{} & \f@tpsg{} & \f@tpsg{} & \f@tpsg{}}

%    \end{macrocode}
% \end{macro}
% \end{macro}
%
% \begin{macro}{\f@tchartstrut}
% \begin{macro}{\ftablewidth}
% \begin{macro}{\fntcolwidth}
% \cs{f@tchartstrut} is a strut used in each table line. \cs{ftablewidth} is 
% width of an externally numbered table. \cs{fntcolwidth} is the minimum
% width of a column in an unnumbered table.
%    \begin{macrocode}
\newcommand*{\f@tchartstrut}{\lower4.5pt\vbox to14pt{}}
\newdimen\ftablewidth
  \ftablewidth=\hsize
\newdimen\fntcolwidth
  \setlength{\fntcolwidth}{0.08\ftablewidth}
%    \end{macrocode}
% \end{macro}
% \end{macro}
% \end{macro}
%
% \begin{macro}{\f@tcol}
% \begin{macro}{\f@tstartchartnonum}
% \cs{f@tstartchartnonum} is a table line of spaces, with no verticals.
%    \begin{macrocode}
\newcommand*{\f@tcol}{%
  \multicolumn{1}{c}{\hspace*{\fntcolwidth}}}
\newcommand*{\f@tstartchartnonum}{%
  \f@tcol &\f@tcol &\f@tcol &\f@tcol &\f@tcol &\f@tcol &\f@tcol &\f@tcol}

%    \end{macrocode}
% \end{macro}
% \end{macro}
%
% \begin{macro}{\ftable}
% \begin{macro}{\f@t@uni@ftable}
% \begin{macro}{\f@tftablenum}
% \begin{macro}{\f@tftablenonum}
% \cs{ftable} sets a complete character table. The actual code is in either
% \cs{f@tftablenum} or \cs{f@tftablenonum} for externally numbered or 
% plain tables, respectively.
% \changes{v1.0a}{2005/12/06}{Added missing zeroing of \cs{f@tn} to \cs{f@tftablenonum}}
%    \begin{macrocode}
\newcommand*{\f@tftablenum}{$$\global\f@tn=\z@
  \halign to\ftablewidth\bgroup
    \f@tchartstrut##\tabskip0pt plus10pt&
    &\hfil##\hfil&\vrule##\cr
    \lower6.5pt\null
    &&&\f@toct0&&\f@toct1&&\f@toct2&&\f@toct3&&\f@toct4&&\f@toct5&&\f@toct6&&\f@toct7&%
    \f@tevenline}
\newcommand*{\f@tftablenonum}{%
  \global\f@tn=\z@
  \begin{tabular}{|c|c|c|c|c|c|c|c|}
    \f@tstartchartnonum
    \f@tevenlinenonum
  \end{tabular}}
\newcommand*{\ftable}{\ifhexoct\f@tftablenum\else\f@tftablenonum\fi}
\newcommand*{\f@t@uni@ftable}{\ifhexoct\f@tftablenum\else\f@tftablenonum\fi}

%    \end{macrocode}
% \end{macro}
% \end{macro}
% \end{macro}
%
% \begin{macro}{\f@tendchart}
% \cs{f@tendchart} sets the last line of an externally numbered table with
% the relevant hex digits.
%    \begin{macrocode}
\newcommand*{\f@tendchart}{\cr\noalign{\hrule}
  \raise11.5pt\null&&&\f@thex 8&&\f@thex 9&&\f@thex A&&\f@thex B&
  &\f@thex C&&\f@thex D&&\f@thex E&&\f@thex F&\cr
  \egroup$$\par}

%    \end{macrocode}
% \end{macro}
%
% \begin{macro}{\f@tpsg}
% \changes{v1.3}{2009/04/30}{Replaced redefinition of \cs{:} by \cs{f@tpsg}}
% \changes{v1.5c}{2009/09/20}{\cs{renewcommand}\cs{f@tpsg} instead (props Michael Niedermair)}
% \changes{v1.5d}{2009/09/22}{Fix the bug for real this time}
% \begin{macro}{\f@placechar}
% \begin{macro}{\f@placedecimal}
% \cs{f@tpsg} typesets a single glyph, possibly with its decimal slot number.
% \cs{f@placechar} is the function to typeset the glyph with its number that is
% internally defined as \cs{f@placedecimal} if decimals are to be shown.
%    \begin{macrocode}
\newcommand*{\f@tpsg}{%
  \setbox\z@=\hbox{\f@placechar{\char\f@tn}{\the\f@tn}}%
  \ifdim\ht\z@>7.5pt\relax
    \f@treposition
  \else
    \ifdim\dp\z@>2.5pt\relax
      \f@treposition
    \fi
  \fi
  \box\z@
  \global\advance\f@tn\@ne
}
%    \end{macrocode}
% Change this definition to adjust the typesetting of the decimal numbers:
%    \begin{macrocode}
\newcommand*\f@placedecimal[2]{#1\ {\tiny #2}}
%    \end{macrocode}
% \end{macro}
% \end{macro}
% \end{macro}
%
% \begin{macro}{\decimals}
% \begin{macro}{\nodecimals}
% Following \cs{decimals}, which is the default, decimal numbers are
% printed in the table. Following \cs{nodecimals} they are not printed.
%    \begin{macrocode}
\newcommand*{\nodecimals}{%
  \renewcommand*\f@placechar{\@firstoftwo}%
}
%    \end{macrocode}
% 
%    \begin{macrocode}
\newcommand{\decimals}{%
  \renewcommand*\f@placechar{\f@placedecimal}%
}
\newcommand*\f@placechar{}
\decimals
%    \end{macrocode}
% \end{macro}
% \end{macro}
%
% \begin{macro}{\f@treposition}
% \cs{f@treposition}
%    \begin{macrocode}
\newcommand*{\f@treposition}{\setbox0=\vbox{\kern2pt\box0}\f@tdim=\dp0
  \advance\f@tdim 2pt \dp0=\f@tdim}

%    \end{macrocode}
% \end{macro}
%
% \begin{macro}{\fonttext}
% \cs{fonttext}\marg{font} typesets \cs{knutext} using \meta{font} (e.g.
% auncl10).
%    \begin{macrocode}
\def\fonttext#1{%
  \def\f@tfontname{#1}%
  \bgroup
  \f@tstartfont
  \knutext
  \egroup}

%    \end{macrocode}
% \end{macro}
%
% \begin{macro}{\regulartext}
% \cs{regulartext}\marg{fontspec} typesets \cs{knutext} using \meta{fontspec} 
% (e.g., \cs{aunclfamily}).
%    \begin{macrocode}
\def\regulartext#1{%
  \bgroup
  #1
  \knutext
  \egroup}

%    \end{macrocode}
% \end{macro}
%
% \begin{macro}{\knutext}
% Deathless prose from Knuth for testing a font. It includes
% \cs{moreknutext}, \cs{capknutext}, and \cs{knunames}.
%    \begin{macrocode}
\def\knutext{{
On November 14, 1885, Senator \& Mrs.~Leland Stanford called together
at their San Francisco mansion the 24~prominent men who had been
chosen as the first trustees of The Leland Stanford Junior University.
They handed to the board the Founding Grant of the University, which
they had executed three days before. This document---with various
amendments, legislative acts, and court decrees---remains as the
University's charter.  In bold, sweeping language it stipulates that
the objectives of the University are ``to qualify students for
personal success and direct usefulness in life; and to promote the
publick welfare by exercising an influence in behalf of humanity and
civilization, teaching the blessings of liberty regulated by law, and
inculcating love and reverence for the great principles of government
as derived from the inalienable rights of man to life, liberty, and
the pursuit of happiness.'' 

\moreknutext

\capknutext

\knunames
\par}}

%    \end{macrocode}
% \end{macro}
%
% \begin{macro}{\@moreknutext}
% Some more text with a variety of ligatures and accents.
%    \begin{macrocode}
\def\@moreknutext{¿But aren't Kafka's Schloß and Æsop's
Œuvres often naïve vis-à-vis the dæmonic phœnix's
official rôle in fluffy soufflés? }

%    \end{macrocode}
% \end{macro}
%
% \begin{macro}{\@capknutext}
% \begin{macro}{\capknutext}
% Text using only capital letters and some punctuation.
%    \begin{macrocode}
\newcommand{\@capknutext}{%
(¡THE DAZED BROWN FOX QUICKLY GAVE 12345--67890 JUMPS!)}
\let\capknutext\@capknutext

%    \end{macrocode}
% \end{macro}
% \end{macro}
%
% \begin{macro}{\@knunames}
% Lots of accents masquerading in personal names.
%    \begin{macrocode}
\def\@knunames{ Ångelå Beatrice Claire
  Diana Érica Françise Ginette Gwdihŵs Hélène Iris
  Jackie Kāaren Łau\.ra María Ni\H{a}tałiĭe Øctave
  Pauline Quêneau Roxanne Sabine Sašo Tã{\'\j}a Tŷ TŶ Uršula
  Vivian Wendy Ŵy Xanthippe Yvønne Zäzilie\par}

%    \end{macrocode}
% \end{macro}
%
% \begin{macro}{\guillemotleft}
% \begin{macro}{\guillemotright}
% \begin{macro}{\flqq}
% \begin{macro}{\frqq}
% Just in case the French quotes are not defined, as they are called for in
% the subsequent \cs{germantext}.
% \changes{v1.2}{2008/05/08}{Replaced \cs{providecommand} for guillemots
%                           by \cs{DeclareTextSymbol}}
%    \begin{macrocode}
\DeclareTextSymbol{\guillemotleft}{OT1}{`\'}
\DeclareTextSymbol{\guillemotright}{OT1}{`\`}
\providecommand{\flqq}{\guillemotleft}
\providecommand{\frqq}{\guillemotright}

%    \end{macrocode}
% \end{macro}
% \end{macro}
% \end{macro}
% \end{macro}
%
% \begin{macro}{\germantext}
% \begin{macro}{\germanparatext}
% Text from the \lpack{Blindtext} package.
% \changes{v1.1}{2006/10/02}{Added \cs{germantext}}
% \changes{v1.51}{2009/05/14}{Changed \cs{germantext} to \cs{germanparatext}}
%    \begin{macrocode}
\providecommand*{\germantext}{%
  \PackageWarning{fonttable}{\protect\germantext\space is deprecated, 
                 \MessageBreak use \protect\germanparatext\space instead}}
\newcommand*{\germanparatext}{%
Dies hier ist ein Blindtext zum Testen von Textausgaben. Wer
diesen Text liest, ist selbst schuld. Der Text gibt lediglich den
Grauwert der Schrift an. Ist das wirklich so? Ist es
gleich\-gül\-tig ob ich schreibe: \frqq Dies ist ein
Blindtext\flqq\ oder \frqq Huardest gefburn\flqq? Kjift --
mitnichten! Ein Blindtext bietet mir wichtige Informationen. An
ihm messe ich die Lesbarkeit einer Schrift, ihre Anmutung, wie
harmonisch die Figuren zueinander stehen und prü\-fe, wie breit
oder schmal sie läuft. Ein Blindtext sollte mög\-lichst viele
verschiedene Buchstaben enthalten und in der Originalsprache
gesetzt sein. Er muß\ keinen Sinn ergeben, sollte aber lesbar
sein. Fremdsprachige Texte wie \frqq Lorem ipsum\flqq\ dienen
nicht dem eigentlichen Zweck, da sie eine
falsche Anmutung vermitteln.\par}

%    \end{macrocode}
% \end{macro}
% \end{macro}
%
% \begin{macro}{\latintext}
% \begin{macro}{\latinparatext}
% The traditional printers' text.
% \changes{v1.1}{2006/10/02}{Added \cs{latintext}}
% \changes{v1.51}{2009/05/14}{Changed \cs{latintext} to \cs{latinparatext} 
% because of a clash with the babel package}
%    \begin{macrocode}
\providecommand*{\latintext}{%
  \PackageWarning{fonttable}{\protect\latintext\space may be overridden by the 
   babel package \MessageBreak use
                 \protect\latinparatext\space instead}}
\newcommand*{\latinparatext}{%
Lorem ipsum dolor sit amet, consectetuer adipiscing elit. Etiam
lobortis facilisis sem. Nullam nec mi et neque pharetra
sollicitudin. Praesent imperdiet mi nec ante. Donec ullamcorper,
felis non sodales commodo, lectus velit ultrices augue, a
dignissim nibh lectus placerat pede. Vivamus nunc nunc, molestie
ut, ultricies vel, semper in, velit. Ut porttitor. Praesent in
sapien. Lorem ipsum dolor sit amet, consectetuer adipiscing elit.
Duis fringilla tristique neque. Sed interdum libero ut metus.
Pellentesque placerat. Nam rutrum augue a leo. Morbi sed elit sit
amet ante lobortis sollicitudin. Praesent blandit blandit mauris.
Praesent lectus tellus, aliquet aliquam, luctus a, egestas a,
turpis. Mauris lacinia lorem sit amet ipsum. Nunc quis urna dictum
turpis accumsan semper.\par}

%    \end{macrocode}
% \end{macro}
% \end{macro}
%
% \begin{macro}{\simpletext}
% \begin{macro}{\fulltext}
% \begin{macro}{\moreknutext}
% \begin{macro}{\knunames}
% \cs{simpletext} kills off \cs{moreknutext} and \cs{knunames}.
% \cs{fulltext} restores \cs{moreknutext} and \cs{knunames}.
% Make \cs{fulltext} the default.
%    \begin{macrocode}
\newcommand*{\simpletext}{\let\moreknutext\relax \let\knunames\relax}
\newcommand*{\fulltext}{\let\moreknutext\@moreknutext \let\knunames\@knunames}
\fulltext

%    \end{macrocode}
% \end{macro}
% \end{macro}
% \end{macro}
% \end{macro}
%
% \begin{macro}{fonttexts}
% \cs{fonttexts}\marg{font}\marg{text} typesets \meta{text} using \meta{font} (e.g.
% auncl10).
% \changes{v1.1}{2006/10/02}{Added \cs{fonttexts}}
%    \begin{macrocode}
\def\fonttexts#1#2{%
  \def\f@tfontname{#1}%
  \bgroup
  \f@tstartfont
  #2
  \egroup}

%    \end{macrocode}
% \end{macro}
%
% \begin{macro}{\regulartexts}
% \cs{regulartext}\marg{fontspec}\marg{text} typesets \meta{text} using \meta{fontspec} 
% (e.g., \cs{aunclfamily}).
% \changes{v1.1}{2006/10/02}{Added \cs{regulartexts}}
%    \begin{macrocode}
\def\regulartexts#1#2{%
  \bgroup
  #1 #2
  \egroup}

%    \end{macrocode}
% \end{macro}
%
% \begin{macro}{\aztext}
% \begin{macro}{\AZtext}
% \begin{macro}{\digitstext}
% \begin{macro}{\punctext}
%  The various characters used for Latin texts.
% \changes{v1.2}{2008/05/08}{Added \cs{aztext}, \cs{AZtext}, \cs{digitstext}
%                            and \cs{punctext}}
%    \begin{macrocode}
\newcommand*{\aztext}{a b c d e f g h i j k l m n o p q r s t u v w x y z}
\newcommand*{\AZtext}{A B C D E F G H I J K L M N O P Q R S T U V W X Y Z}
\newcommand*{\digitstext}{0 1 2 3 4 5 6 7 8 9}
\newcommand*{\punctext}{` ! @ \$ \& * ( ) \_ - + =  [ ] < > \{ \} : ; ' , . ? /}

%    \end{macrocode}
% \end{macro}
% \end{macro}
% \end{macro}
% \end{macro}
%
% \subsection{Testing a glyph}
%
% This is a reimplementation of Donald Knuth's \texttt{testfont.tex} which
% is available from CTAN and there is also a commented version in Appendix H
% of \textit{The METAFONT Book}.
%
% \begin{macro}{\fnthours}
% \begin{macro}{\f@ttwodigits}
% The time of day on a 24 hour clock.
%    \begin{macrocode}
%%%%%%%% using \@tempcnta for Knuth's \m and \@tempcntb for his \n
\newcommand*{\fnthours}{\@tempcntb=\time \divide\@tempcntb 60
  \@tempcnta=-\@tempcntb \multiply\@tempcnta 60 \advance\@tempcnta \time
  \f@ttwodigits\@tempcntb:\f@ttwodigits\@tempcnta}
\newcommand*{\f@ttwodigits}[1]{\ifnum #1<10 0\fi \number#1}

%    \end{macrocode}
% \end{macro}
% \end{macro}
%
% \begin{macro}{\f@tgettsechars}
% \begin{macro}{\f@ttchar}
% \begin{macro}{\f@tschar}
% \begin{macro}{\f@techar}
% \cs{f@tgettsechars}\marg{T}\marg{S}\marg{E} gets three characters
% and \cs{chardef}s \cs{f@ttchar} to \meta{T} (the test character),
% \cs{f@tschar} to \meta{S} (start character) and \cs{f@techar}
% to \meta{E} (the end character).
%    \begin{macrocode}
\newcommand*{\f@tgettsechars}[3]{%
  \chardef\f@ttchar=#1 \chardef\f@tschar=#2 \chardef\f@techar=#3}

%    \end{macrocode}
% \end{macro}
% \end{macro}
% \end{macro}
% \end{macro}
%
% \begin{macro}{\glyphmixture}
% \begin{macro}{\f@tmixpattern}
% \begin{macro}{\f@tdomix}
%  \cs{glyphmixture}\marg{T}\marg{S}\marg{E} sets a mix of \meta{T} within
% the glyph range from \meta{S} to \meta{E} according to the pattern
% \cs{f@tmixpattern}. The work is done by \cs{f@tdomix}.
%    \begin{macrocode}
\newcommand*{\glyphmixture}[3]{\f@tgettsechars{#1}{#2}{#3}%
                               \f@tdomix\f@tmixpattern}
\newcommand*{\f@tmixpattern}{\f@tocta\f@toctb\f@tocta\f@tocta\f@toctb\f@toctb\f@tocta\f@tocta\f@tocta\f@toctb\f@toctb\f@toctb\f@tocta\f@toctb}
\newcommand*{\f@tdomix}[1]{\par\chardef\f@tocta=\f@ttchar \@tempcntb=\f@tschar
   \loop \chardef\f@toctb=\@tempcntb #1\endgraf
   \ifnum \@tempcntb<\f@techar \advance\@tempcntb \@ne \repeat}

%    \end{macrocode}
% \end{macro}
% \end{macro}
% \end{macro}
%
% \begin{macro}{\glyphalternation}
% \begin{macro}{\f@taltpattern}
% These are similar to \cs{glyphmixture} and \cs{f@tmixpattern} except
% that the glyphs are alternated.
%    \begin{macrocode}
\newcommand*{\glyphalternation}[3]{\f@tgettsechars{#1}{#2}{#3}%
                                   \f@tdomix\f@taltpattern}
\newcommand*{\f@taltpattern}{\f@tocta\f@toctb\f@tocta\f@toctb\f@tocta\f@toctb\f@tocta\f@toctb\f@tocta\f@toctb\f@tocta\f@toctb\f@tocta\f@toctb\f@tocta\f@toctb\f@tocta}

%    \end{macrocode}
% \end{macro}
% \end{macro}
%
% \begin{macro}{\f@tdisc}
% For breaking long lines so that the test character will be at the end
% of one line and repeated at the start of the next one.
%    \begin{macrocode}
\newcommand*{\f@tdisc}{\discretionary{\f@ttchar}{\f@ttchar}{\f@ttchar}}

%    \end{macrocode}
% \end{macro}
%
% \begin{macro}{\glyphseries}
% \begin{macro}{\f@tdoseries}
% \cs{glyphseries}\marg{T}\marg{S}\marg{E} puts the test character \meta{T}
% between all the others in the range \meta{S} to \meta{E}. The work is
% done by \cs{f@tdoseries}.
%    \begin{macrocode}
\newcommand*{\glyphseries}[3]{\f@tgettsechars{#1}{#2}{#3}%
  \f@tdisc\f@tdoseries\f@tschar\f@techar\par}
\newcommand*{\f@tdoseries}[2]{\@tempcntb=#1\relax
  \loop\char\@tempcntb\f@tdisc
    \ifnum\@tempcntb<#2\advance\@tempcntb \@ne \repeat}

%    \end{macrocode}
% \end{macro}
% \end{macro}
%
% \begin{macro}{\glyphalphabet}
% \begin{macro}{\GLYPHALPHABET}
% \begin{macro}{\f@tcomplower}
% \begin{macro}{\f@tcompupper}
% \cs{glyphalphabet}\marg{T} inserts the test glyph \meta{T} between the lowercase
% alphabetic characters. Similarly \cs{GLYPHALPHABET}\marg{T} does
% the same with the uppercase characters. The work is done by, respectively,
% \cs{f@tcomplower} and \cs{f@tcompupper}.
%    \begin{macrocode}
\newcommand*{\glyphalphabet}{\f@tcomplower}
\newcommand*{\GLYPHALPHABET}{\f@tcompupper}
\newcommand*{\f@tcomplower}[1]{\chardef\f@ttchar=#1 
  \f@tdisc\f@tdoseries{`a}{`z}\f@tdoseries{31}{34}\par}
\newcommand*{\f@tcompupper}[1]{\chardef\f@ttchar=#1 
  \f@tdisc\f@tdoseries{`A}{`Z}\f@tdoseries{35}{37}\par}

%    \end{macrocode}
% \end{macro}
% \end{macro}
% \end{macro}
% \end{macro}
%
% \begin{macro}{\glyphlowers}
% \begin{macro}{\glyphuppers}
% \begin{macro}{\glyphdigits}
% \begin{macro}{\f@tclc}
% \begin{macro}{\f@tcuc}
% \begin{macro}{\f@tdgs}
% \begin{macro}{\f@tdocomprehensive}
% These macros generate an extended mix of characters of a particular kind.
% The work is done by \cs{f@tdocomprensive} with \cs{f@tclc}, \cs{f@tcuc}, and
% \cs{f@tdgs} setting up the glyph sets.
%    \begin{macrocode}
\newcommand*{\glyphlowers}{\f@tdocomprehensive\f@tclc{`a}{`z}{31}{34}}
\newcommand*{\glyphuppers}{\f@tdocomprehensive\f@tcuc{`A}{`Z}{35}{37}}
\newcommand*{\glyphdigits}{\f@tdocomprehensive\f@tdgs{`0}{`4}{`5}{`9}}
\newcommand*{\f@tdocomprehensive}[5]{\par\chardef\f@ttchar=#2
  \loop{#1} \ifnum\f@ttchar<#3\@tempcnta=\f@ttchar\advance\@tempcnta \@ne
  \chardef\f@ttchar=\@tempcnta \repeat
  \chardef\f@ttchar=#4
  \loop{#1} \ifnum\f@ttchar<#5\@tempcnta=\f@ttchar\advance\@tempcnta \@ne
  \chardef\f@ttchar=\@tempcnta \repeat}
\newcommand*{\f@tclc}{\f@tdisc\f@tdoseries{`a}{`z}\f@tdoseries{31}{34}\par}
\newcommand*{\f@tcuc}{\f@tdisc\f@tdoseries{`A}{`Z}\f@tdoseries{35}{37}\par}
\newcommand*{\f@tdgs}{\f@tdisc\f@tdoseries{`0}{`9}\par}

%    \end{macrocode}
% \end{macro}
% \end{macro}
% \end{macro}
% \end{macro}
% \end{macro}
% \end{macro}
% \end{macro}
%
% \begin{macro}{\glyphpunct}
% \begin{macro}{\f@tdopunct}
% \cs{glyphpunct} sets punctuation marks in combination with different
% sorts of letters. The work is done by \cs{f@tdopunct}.
%    \begin{macrocode}
\newcommand*{\glyphpunct}{\par\f@tdopunct{min}\f@tdopunct{pig}\f@tdopunct{hid}
                     \f@tdopunct{HIE}\f@tdopunct{TIP}\f@tdopunct{fluff}
  \$1,234.56 + 7/8 = 9\% @ \#0\par}
\newcommand*{\f@tdopunct}[1]{#1,\ #1:\ #1;\ ‘#1’\
  ¿#1?\ ¡#1!\ (#1)\ [#1]\ #1*\ #1.\par}

%    \end{macrocode}
% \end{macro}
% \end{macro}
%
% The end of the package.
%    \begin{macrocode}
%</sty>
%    \end{macrocode}
%
%    \begin{macrocode}
%<*package>
%    \end{macrocode}
%
%    By default the package uses coloring to improve the table
%    appearance and therefore requires a color package.
%    \begin{macrocode}
\RequirePackage{xcolor}
%    \end{macrocode}
%

%    \begin{macrocode}
%<@@=fnttable>
%    \end{macrocode}
%
%    We need the package \pkg{xparse} for specifying the document-level
%    interface commands and \pkg{l3keys2e} to use the \pkg{expl3} key
%    value methods within \LaTeXe{}. These packages automatically
%    require \pkg{expl3} so there is no need to load that explicitly.
%    Actually, \pkg{expl3}, \pkg{l3keys2e} and the \pkg{xparse}
%    functionality is now all part of the \LaTeX{} kernel so the next
%    line is actually not needed at all with a current \LaTeX{} kernel, but
%    in order to support older installations we keep it for now.
%
%    \begin{macrocode}
\RequirePackage{xparse,l3keys2e}
%    \end{macrocode}
%
%  \newcommand\hex[1]{$\langle\textit{hex}_{#1}\rangle$}
%
%                    
% \subsection{User interface commands}
%
%
%  Throughout the implementation we will define a number of keys (and
%  their allowed values). We introduce them at the point where they are
%    used, so they are sprinkled throughout the code.\footnote{This fits
%    with the way this package was developed. I first implemented a
%    single rigid table layout without configuration possibilities and
%    then thought about which parts I wanted to have flexible. I then
%    replaced the rigid code with code that is affected by setting
%    key/value pairs.}

%  \begin{macro}{\fonttablesetup}
%    To set up user defaults for the keys we provide a standard
%    interface. The command \cs{fnttable} expects a
%    key/value list and can be called as often as necessary.
%    \begin{macrocode}
\ExplSyntaxOn
\NewDocumentCommand \fonttablesetup { m }
  { \keys_set:nn {@@} {#1} \ignorespaces }
%    \end{macrocode}
%  \end{macro}
%
%
%
%  \begin{macro}{\displayfonttable}
%    The document-level command for generating a font table.
%    \begin{macrocode}
\NewDocumentCommand\displayfonttable {O{} mmmm}{%
%    \end{macrocode}
%    For the starred form we preset a number of keys with values
%    suitable when displaying 8-bit legacy fonts.
%    With such fonts Unicode block headers make little
%    sense (as the fonts do not conform to the Unicode layout and
%    since they have at most 265 glyphs). It is therefore also unnecessary to
%    loop over the whole Unicode range of the first plane.
%    If necessary all of them can still be overwritten in the optional
%    argument.
%       
%    Some font names contain special character, e.g., \verb=_= or
%    \verb=&=.   We therefore turn the font name into a string before
%    passing it on but we allow for the possibility that the font name
%    was given in a macro, so we do this with \cs{tl_to_str:o} (tough
%    if the font name starts out with a special character then
%    you have to use \cs{string} or something to prevent the special
%    character from causing issues).
% \changes{v1.0k}{2025/07/11}{Allow special chars in font names but
%    not as the first character (gh/12)}
%    \begin{macrocode}
  \str_if_case:enTF {#2} {
    {TU} {}
    {tu} {}
    {EU1} {}
    {eu1} {}
    {EU2} {}
    {eu2} {}
  } {
    \@@_display_fonttable:nnnnn {#2}{#3}{#4}{#5}{#6}
  } {
    \@@_display_fonttable:nnnnn
    {nostatistics,display-block=none,hex-digits=head,range-end=FF,#2}
    {#3}{#4}{#5}{#6}
  }
}
%    \end{macrocode}
%  \end{macro}
%  
%  
%  \begin{macro}{\@@_display_fonttable:nnnnn}
%    This command is the main workhorse of the
%    package. It produces a \texttt{longtable} containing all font
%    glyphs with 16 glyphs per row. The first optional argument is
%    used to configure the table through key/value pairs, the
%    mandatory argument is the font name to display (in
%    \texttt{fontspec} conventions) and the final optional argument is
%    the font feature list if any. If the latter is not provided it
%    will get a special value (\texttt{--NoValue--}) assigned by
%    \texttt{xparse}, which is something that can be tested for.
%
%    First: optional keys; second through fifth: fontspec.
%    \begin{macrocode}
\cs_new:Npn \@@_display_fonttable:nnnnn #1#2#3#4#5 {
  \group_begin:
%    \end{macrocode}
%    First initialize the font that should be displayed (perhaps with a
%    feature list) and then update the key/value list using \verb=#1=. 
%    \begin{macrocode}
    \fontencoding{#2}\fontfamily{#3}\fontseries{#4}\fontshape{#5}
    \selectfont
    \keys_set:nn{@@}{#1}
%    \end{macrocode}
%    If the \LuaTeX{} engine is used without HarfBuzz and the display
%    range includes code points above \texttt{U+EFFFF} the output shows
%    remappings and not what is in the font, so we issue a warning.
% \changes{v1.0g}{2022/11/12}{Test for luatex without harfbuzz and
%    private area A and warn (gh/8)}
%    \begin{macrocode}
    \bool_lazy_and:nnT
        { \sys_if_engine_luatex_p: }
        { \int_compare_p:nNn { "EFFFF } < { "\l_@@_range_end_tl } }
        {
          \directlua{token.put_next(token.create(font.getfont(font.current()).hb~
                     and~ 'use_none:n'~  or~ 'use:n'))}
            { \msg_warning:nn {fnttable}{noharfbuzz} }
        }
%    \end{macrocode}
%    If the user has asked for a comparsion to some other font we need to set this up:
%    \begin{macrocode}
    \tl_if_empty:NTF \l_@@_compare_with_tl
       { \tl_clear:N \l_@@_compare_font_tl }
       {
         \setfontface \l_@@_compare_font_tl {\l_@@_compare_with_tl}[]
         \cs_set_eq:NN \@@_handle_missing_glyph:n
                       \@@_handle_missing_glyph_compare:n
       }
%    \end{macrocode}
%    Typesetting the font tables in twocolumn mode makes little sense
%    due to their width, and if \env{longtable} is used it will
%    complain. However there is one case where it should work: in a
%    page-wide float. To make this happen we claim that we are not in
%    twocolumn mode if the display is inside a vertical box.
% \changes{v1.0g}{2022/11/12}{Support use in twocolumn mode if inside
%    a table* float (gh/7)}
%    \begin{macrocode}
    \if_mode_vertical: \if_mode_inner: \@twocolumnfalse \fi: \fi:
%    \end{macrocode}
%    Then we start the table with 17 columns. We use \texttt{longtable}
%    if we produce a caption and \texttt{longtable*} if not (so that
%    the table number is not increased, which would look odd if you
%    have other tables in your document).
%    \begin{macrocode}
    \begin{longtable\bool_if:NF\l_@@_display_header_bool{*}}
          {@{}r@{\quad}*{16}{c}@{}}
%    \end{macrocode}
%    Special headers and footers are set up first:
%    \begin{macrocode}
      \@@_setup_header_footer:nn{#2}{#3}
%    \end{macrocode}
%    Then we produce all table rows with the glyphs.
%    \begin{macrocode}
      \@@_produce_table_rows:
%    \end{macrocode}
%    At the very end we may typeset some statistics. This can't be
%    done in the table footer, because the data is dynamic (e.g.,
%    number of glyphs processed) and the table footers are static and
%    do not change based on the table content.
%    \begin{macrocode}
      \@@_handle_table_ending:n {#2}
    \end{longtable\bool_if:NF\l_@@_display_header_bool{*}}
  \group_end:
}
%    \end{macrocode}
%  \end{macro}
%
%
%    \begin{macrocode}
\msg_new:nnn {fnttable}{noharfbuzz}
    { You~ asked~ for~ displaying~ glyphs~ with~ code \iow_newline:
      points~ above~ U+EFFFF~ \msg_line_context: ,~ i.e.,~ from~ the~
      'Supplementary~ Private~ Use~ Area-A'\iow_newline:
      without~ specifying~ '[Renderer=Harfbuzz]'~ when~
      loading~ the~ font.
      \iow_newline:\iow_newline:
      With~ LuaLaTeX,~ this~ Unicode~ region~ is~ used~
      for~ remappings~ (if~ the~ HarfBuzz~ engine~ is~ not~ used).~
      Thus,~ the~ results~ shown~ do~ not~ reflect~ what~
      is~ in~ the~ font!
    }
%    \end{macrocode}
%
%
%
%  \begin{macro}{\fonttableglyphcount}
%  \begin{macro}{\g_@@_glyph_int,\g_@@_glyph_only_B_int,\g_@@_glyph_also_B_int}
%    While generating the font table we count the number of glyphs we
%    see (and typeset). The total is available in the command
%    \cs{fonttableglyphcount} after the table got finished and will be reset to
%    zero when the next table starts.
%    \begin{macrocode}
\DeclareDocumentCommand \fonttableglyphcount {}
                        { \int_use:N \g_@@_glyph_int }
%    \end{macrocode}
%    
%    \begin{macrocode}
\int_new:N \g_@@_glyph_int
%    \end{macrocode}
%    When comparing fonts we also record data for the second font: the
%    number of glyphs in both and the number of glyphs only in the
%    second one.
%    \begin{macrocode}
\int_new:N \g_@@_glyph_only_B_int
\int_new:N \g_@@_glyph_also_B_int
%    \end{macrocode}
%  \end{macro}
%  \end{macro}




% \subsection{The overall table layout}
%
%     
%  \begin{macro}{\@@_setup_header_footer:nn}
%    Setting up header and footer lines of the table.
%    This macro receives the \textit{font name} and the \textit{font
%    features} specified by the user as its arguments.
%    \begin{macrocode}
\cs_new:Npn \@@_setup_header_footer:nn #1#2{
%    \end{macrocode}
%    On the first page of the table the header may show a caption or
%    some other sort of title based on the value of
%    \cs{l_@@_display_header_bool}. The formatting is handled by
%    \cs{@@_format_table_title:nn} which can be customized through the
%    key \key{title-format}.
%    \begin{macrocode}
    \bool_if:NT \l_@@_display_header_bool
       { \@@_format_table_title:nn{#1}{#2} \@@_debug_nl:n{T}\\*[6pt] }
%    \end{macrocode}
%    We may also want to display a line of hex digits. This is
%    controlled through the key \key{hex-digits} that accepts different
%    values: \kval{head}, \kval{foot}, \kval{head+foo}, \kval{block}
%    (after a block title) or \kval{none}.
%    \begin{macrocode}
    \bool_if:NT \l_@@_header_hex_digits_bool
      { \@@_display_row_of_hex_digits:   \@@_debug_nl:n{H}\\*      }
  \endfirsthead
%    \end{macrocode}
%    Headers for later table pages have a continuation title and
%    maybe a row of hex digits.
%    \begin{macrocode}
    \bool_if:NT \l_@@_display_header_bool
      { \@@_format_table_cont:nn{#1}{#2} \@@_debug_nl:n{T}\\*[6pt] }
    \bool_if:NT \l_@@_header_hex_digits_bool
      { \@@_display_row_of_hex_digits:   \@@_debug_nl:n{H}\\*      }
  \endhead
%    \end{macrocode}
%    Footers of the table are either empty or show a row of hex digits.
%    \begin{macrocode}
    \bool_if:NT \l_@@_footer_hex_digits_bool
      { \@@_display_row_of_hex_digits:   \@@_debug_nl:n{H}\\*      }
  \endfoot
%    \end{macrocode}
%    The footer of the last page of the table will always be
%    empty. Any special row, such as a row of hex digits, will be
%    provided in the table body. The reason is that we may want to
%    display statistics at the very end of the table and those can't be
%    placed into a static footer.
%    \begin{macrocode}
  \endlastfoot
}
%    \end{macrocode}
% \end{macro}
%
%  \begin{macro}{\l_@@_header_hex_digits_bool}
%  \begin{macro}{\l_@@_footer_hex_digits_bool}
%  \begin{macro}{\l_@@_blockwise_hex_digits_bool}
%    Here are the booleans we use in the code.
%    \begin{macrocode}
\bool_new:N \l_@@_header_hex_digits_bool     
\bool_new:N \l_@@_footer_hex_digits_bool
\bool_new:N \l_@@_blockwise_hex_digits_bool
%    \end{macrocode}
%  \end{macro}
%  \end{macro}
%  \end{macro}
%    
%    
%  \begin{macro}{\@@_display_row_of_hex_digits:}
%  \begin{macro}{\@@_format_hex_digit:n}
%    Producing a row of hex digits is simple.
%    \begin{macrocode}
\cs_new:Npn \@@_display_row_of_hex_digits: {
    & \@@_format_hex_digit:n{0}     & \@@_format_hex_digit:n{1} 
    & \@@_format_hex_digit:n{2}     & \@@_format_hex_digit:n{3} 
    & \@@_format_hex_digit:n{4}     & \@@_format_hex_digit:n{5} 
    & \@@_format_hex_digit:n{6}     & \@@_format_hex_digit:n{7} 
    & \@@_format_hex_digit:n{8}     & \@@_format_hex_digit:n{9} 
    & \@@_format_hex_digit:n{A}     & \@@_format_hex_digit:n{B} 
    & \@@_format_hex_digit:n{C}     & \@@_format_hex_digit:n{D} 
    & \@@_format_hex_digit:n{E}     & \@@_format_hex_digit:n{F}  }
%    \end{macrocode}
%    Each digit is typeset in typewriter and in script size. We offer
%    font and color
%    customizations. Note that it is important to set an explicit
%    family. Otherwise the hex digits are formatted using the current
%    table font (which may or may not work at all).
%    \begin{macrocode}
\cs_new:Npn \@@_format_hex_digit:n #1 {
  \l_@@_hex_digits_font_tl \l_@@_color_tl #1 }
%    \end{macrocode}
%  \end{macro}
%  \end{macro}
%

%  \begin{macro}{\l_@@_color_tl}
%    The token list to hold definition if set up.
%    \begin{macrocode}
\tl_new:N \l_@@_color_tl
%    \end{macrocode}
%  \end{macro}
%
%    
% \keysetup{overall table}
%    Here are the definitions for the keys used in the code above:
%    \begin{macrocode}
\keys_define:nn {@@} {
%    \end{macrocode}
%    The \key{header} key is a boolean that determines if a header
%    title should be produced (default)
%    \begin{macrocode}
  ,header .bool_set:N   = \l_@@_display_header_bool
  ,header .default:n    = true
  ,header .initial:n    = true
%    \end{macrocode}
%    To ease the setup we also support the key \key{noheader} which is
%    a short form for \texttt{header=false}.
%    \begin{macrocode}
  ,noheader .bool_set_inverse:N = \l_@@_display_header_bool
  ,noheader .default:n          = true
%    \end{macrocode}
%    The default for the \key{title-format} key is to produce a
%    \cs{caption} listing the font name and any features (if
%    given). Note the \cs{IfValueTF} command (provided by
%    \texttt{xparse}) that checks if the second argument got any value
%    or has the special \texttt{--NoValue--} value.
%    \begin{macrocode}
  ,title-format      .cs_set:Np = \@@_format_table_title:nn #1#2
  ,title-format      .initial:n = 
  \IfValueTF{#2} { \caption{ #1~ (features:~ \texttt{\small#2}) } } 
      { \caption{ #1 } }
%    \end{macrocode}
%    The default continuation title ignores the given features, so the
%    formatting is somewhat simpler. It uses \verb=\caption[]{...}= to
%    make a caption that doesn't alter the table number.
% \changes{v1.0f}{2021/10/29}{Make documentation and code match:
%                             it should be \texttt{title-format-cont}}
%    \begin{macrocode}
  ,title-format-cont .cs_set:Np = \@@_format_table_cont:nn #1#2
  ,title-format-cont .initial:n = \caption[]{#1~ \emph{cont.}}
%    \end{macrocode}
%    The key \key{hex-digits} is implemented as a choice, where each
%    allowed value sets different booleans that are then used in the code.
%    \begin{macrocode}
  ,hex-digits  .choice:
  ,hex-digits / block   .code:n  =
  \bool_set_true:N   \l_@@_blockwise_hex_digits_bool
  \bool_set_false:N  \l_@@_header_hex_digits_bool
  \bool_set_false:N  \l_@@_footer_hex_digits_bool
  ,hex-digits / foot   .code:n  =
  \bool_set_true:N   \l_@@_footer_hex_digits_bool
  \bool_set_false:N  \l_@@_header_hex_digits_bool
  \bool_set_false:N  \l_@@_blockwise_hex_digits_bool
  ,hex-digits / head   .code:n  =
  \bool_set_true:N   \l_@@_header_hex_digits_bool
  \bool_set_false:N  \l_@@_footer_hex_digits_bool
  \bool_set_false:N  \l_@@_blockwise_hex_digits_bool
  ,hex-digits / head+foot  .code:n  =
  \bool_set_true:N   \l_@@_header_hex_digits_bool
  \bool_set_true:N   \l_@@_footer_hex_digits_bool
  \bool_set_false:N  \l_@@_blockwise_hex_digits_bool
  ,hex-digits / none   .code:n  =
  \bool_set_false:N  \l_@@_header_hex_digits_bool
  \bool_set_false:N  \l_@@_footer_hex_digits_bool
  \bool_set_false:N  \l_@@_blockwise_hex_digits_bool
  ,hex-digits  .initial:n = head
%    \end{macrocode}
%    The font for hex digits are set with \key{hex-digits-font}.
%    \begin{macrocode}
  ,hex-digits-font  .tl_set:N  = \l_@@_hex_digits_font_tl
  ,hex-digits-font  .initial:n = \ttfamily \scriptsize
%    \end{macrocode}
%    Customizing the row header (on the left) can be done with this
%    key. Defaults for font, fontsize, and color is set on the outside, but can, of
%    course, be overwritten inside if that is desired.
% \changes{v1.0g}{2022/11/12}{Add key hex-digits-row-format to allow
%    customizing the row title on the left (gh/3)}
%    \begin{macrocode}
  ,hex-digits-row-format  .cs_set:Np = \@@_format_row_hex_digits:n #1
  ,hex-digits-row-format  .initial:n = U+#1 0 \, - \, #1 F
%    \end{macrocode}
%    The \key{color} key is used in  most places that get colored; some
%    have their own key but default to the main color.         
%    \begin{macrocode}
  ,color .choice:
  ,color / none    .code:n    = \tl_clear:N \l_@@_color_tl
  ,color / unknown .code:n    = \tl_set:Nn \l_@@_color_tl { \color {#1} }
  ,color           .initial:n = blue 
}
%    \end{macrocode}
%
%
%
%
%  \begin{macro}{\@@_handle_table_ending:n}
%    At the end of the table we may want to display a final row of
%    hex digits and perhaps some statistics, i.e., the number of
%    typeset glyphs.
%    \begin{macrocode}
\cs_new:Npn \@@_handle_table_ending:n #1 {
%    \end{macrocode}
%     
%    \begin{macrocode}
  \@@_debug_nl:n{H} \\*
  \bool_if:NT \l_@@_footer_hex_digits_bool
    { \@@_display_row_of_hex_digits: \@@_debug_nl:n{H} \\*   }
  \bool_if:NT \l_@@_display_statistics_bool
    { \\*[2pt]
      \multicolumn{17}{l}{ \l_@@_stats_font_tl
%    \end{macrocode}
%    If we do font comparison, we use a different command for
%    displaying statistics and pass more data to it.
%    \begin{macrocode}
        \tl_if_empty:NTF \l_@@_compare_with_tl {
          \@@_format_stats:nn{#1}{\fonttableglyphcount}
        } {
          \@@_format_compare_stats:nnnnnn{#1}{\fonttableglyphcount}
          { \l_@@_compare_with_tl }
%    \end{macrocode}
%    The extra arguments are total glyph number in second font, glyphs
%    missing in second font and glyphs only in second font.
%    \begin{macrocode}
          { \int_eval:n { \int_use:N\g_@@_glyph_also_B_int +
            \int_use:N\g_@@_glyph_only_B_int }
          }
          { \int_eval:n { \fonttableglyphcount -
            \int_use:N\g_@@_glyph_also_B_int }
          }
          { \int_use:N\g_@@_glyph_only_B_int }
        }
%    \end{macrocode}
%    We don't know exactly how wide the table is (and nor does the
%    user) but one may need to use \cs{parbox} when formatting
%    the statistic line(s). So we back up a bit (rather random) which
%    allows us to use \verb=\parbox{\linewidth}{...= in the key
%    without thinking too much about it.
%    \begin{macrocode}
        \hspace*{-3cm}
      }
    }
}
%    \end{macrocode}
%
% \keysetup{for statistics}
%    Here are the keys used above. By default we produce statistics.
%    \begin{macrocode}
\keys_define:nn {@@} {
    ,statistics .bool_set:N = \l_@@_display_statistics_bool
    ,statistics .default:n   = true
    ,statistics .initial:n   = true
%    \end{macrocode}
%    the key \key{nostatistics} is just short for \texttt{statistics=false}:
%    \begin{macrocode}
    ,nostatistics .bool_set_inverse:N = \l_@@_display_statistics_bool
    ,nostatistics .default:n = true
%    \end{macrocode}
%   The default font we use is \cs{normalfont}. Again we need to
%    supply a family to avoid getting the font used in the table body. 
%    \begin{macrocode}
    ,statistics-font    .tl_set:N  = \l_@@_stats_font_tl
    ,statistics-font    .initial:n = \normalfont\small
%    \end{macrocode}
%    And here we have the default text. There is only space for a
%    single line. If more text is needed one needs to provide some
%    explicit \cs{parbox}.
% \changes{v1.0g}{2022/11/12}{Change default text so that it makes
%    more sense if only a portion of the font is displayed (gh/4)}
%    \begin{macrocode}
    ,statistics-format  .cs_set:Np = \@@_format_stats:nn #1#2
    ,statistics-format  .initial:n = Total~ number~ of~ glyphs~ shown~ from~ #1:~#2 
  }
%    \end{macrocode}
%  \end{macro}
%
%
%
%
%
%  \begin{macro}{\@@_debug_nl:n}
%    While developing the code I had a bit of trouble getting the line
%    endings correct, so I added a little macro that made them visible
%    (displaying its argument in the table margin when the key
%    \key{debug} is used. By default it does nothing.
%    \begin{macrocode}
\cs_new:Npn \@@_debug_nl:n #1 {}
%    \end{macrocode}
%    
% \keysetup{debugging}
%    This key is really internal and is therefore not documented above
%    (and its behavior may changes over time).
%    \begin{macrocode}
\keys_define:nn {@@} {
  debug .code:n = \cs_set:Npn \@@_debug_nl:n ##1
      {\rlap{\normalfont\scriptsize \qquad ##1}} }
%    \end{macrocode}
%  \end{macro}
%    
%
%
%
%
%
%
% \subsection{The producing the table content}
%
%
%    The body of the table consists of rows with sixteen glyphs each
%    and to produce it we loop through all possible Unicode points
%    starting at \texttt{U+0000} and ending with \texttt{U+FFFF}.
%
%    This is implemented with a four-level nested loop that runs through the
%    values \texttt{0}, \texttt{1}, \ldots, \texttt{F} with the
%    current hex value in each of the four positions stored in some variable.


%  \begin{macro}{\g_@@_hex_H_tl,\g_@@_hex_A_tl,
%                \g_@@_hex_B_tl,\g_@@_hex_C_tl}
%    \cs{g_@@_hex_H_tl} is a bit special because, it is initially not
%    zero, but empty, so that slots in the lower plane are denoted by 4
%    hex digits.
%    We really only need three further variables, as the value in the
%    innermost loop can used directly.
%    \begin{macrocode}
\tl_new:N \g_@@_hex_H_tl   % higher plane
\tl_new:N \g_@@_hex_A_tl
\tl_new:N \g_@@_hex_B_tl
\tl_new:N \g_@@_hex_C_tl
%    \end{macrocode}
%  \end{macro}
%
%
%  \begin{macro}{\c_@@_hex_digits_clist}
%    Here is the sequence we loop through on each level, except the
%    one for the outer level.
%    \begin{macrocode}
\clist_const:Nn\c_@@_hex_digits_clist{0,1,2,3,4,5,6,7,8,9,A,B,C,D,E,F}
%    \end{macrocode}
%  \end{macro}
%
%
%
%  \begin{macro}{\@@_produce_table_rows:,\@@_handle_hex_H:n,
%                \@@_handle_hex_A:n,\@@_handle_hex_B:n,
%                \@@_handle_hex_C:n,\@@_handle_hex_D:n}
%    The overall code layout is then fairly simply:
%    \begin{macrocode}
\cs_new:Npn \@@_produce_table_rows: {
%    \end{macrocode}
%    First to some general initialization
%    \begin{macrocode}
  \@@_initialize_table_rows:
%    \end{macrocode}
%    and then loop we start the loop. The outer level is a bit special
%    as currently Unicode has only slots allocated in plane 0, 1, 2
%    and E (well, and F, but that is a private area) so we loop only
%    over those and instead of \texttt{0} we use an empty value.
%    Not covered is the whole of plane 16 which too is now a
%    private area.  
%
%    \begin{macrocode}
   \clist_map_function:nN { { } , 1, 2, E, F } \@@_handle_hex_H:n }
%    \end{macrocode}
%    
%    Most fonts do not have glyphs in the higher planes, which is why
%    by default we don't loop using a nonempty \cs{@@_handle_hex_H:n}.
%    But if the user wants to scan and display the higher slots they
%    can by setting \key{range-end} appropriatly.
%
%    So after setting \cs{@@_handle_hex_H:n} we loop over
%    \cs{c_@@_hex_digits_clist} for the next
%    hex digit (which we call \enquote{A}).
%    \begin{macrocode}
\cs_new:Npn \@@_handle_hex_H:n #1 { \tl_gset:Nn\g_@@_hex_H_tl{#1}
   \clist_map_function:NN \c_@@_hex_digits_clist \@@_handle_hex_A:n }
%    \end{macrocode}
%
%    Handling \enquote{A} means storing its value for later use and
%    then start a loop for setting the second (or third on higher planes) hex digits:
%    \begin{macrocode}
\cs_new:Npn \@@_handle_hex_A:n #1 { \tl_gset:Nn\g_@@_hex_A_tl{#1}
   \clist_map_function:NN \c_@@_hex_digits_clist \@@_handle_hex_B:n }
%    \end{macrocode}
%
%    Same game for \enquote{B} and \enquote{C}\footnote{Actually this
%    is a white lie. In reality we do a lot of extra stuff when
%    handling \enquote{C} so later one we give a second definition for
%    \cs{@@_handle_hex_C:n} but for understanding the overall picture
%    the simpler one shown here is better.}:
%    \begin{macrocode}
\cs_new:Npn \@@_handle_hex_B:n #1 { \tl_gset:Nn\g_@@_hex_B_tl{#1}
   \clist_map_function:NN \c_@@_hex_digits_clist \@@_handle_hex_C:n }
%    \end{macrocode}
%
%    \begin{macrocode}
\cs_new:Npn \@@_handle_hex_C:n #1 { \tl_gset:Nn\g_@@_hex_C_tl{#1}
   \clist_map_function:NN \c_@@_hex_digits_clist \@@_handle_hex_D:n }
%    \end{macrocode}
%    In the innermost loop we now have the full Unicode number
%    available, so there we have to decide what to do with it. This is
%    done by \cs{@@_handle_hex_D:n} that receives the full number,
%    e.g., \texttt{1A7C} or \texttt{1AD00}, as its argument.
%    \begin{macrocode}
\cs_new:Npn \@@_handle_hex_D:n #1 {
  \@@_handle_slot:x
     { " \g_@@_hex_H_tl \g_@@_hex_A_tl
         \g_@@_hex_B_tl \g_@@_hex_C_tl #1 }
}
%    \end{macrocode}
%  \end{macro}
%
%  \begin{macro}{\g_@@_row_tl}
%    We first collect the glyphs for a whole row before deciding to
%    typeset it, because if the row is entirely empty we want to omit
%    it. The data for the row is collected slot by slot and the typesetting
%    information (the glyph or the indication for a missing glyph is
%    appended to \cs{g_@@_row_tl}.
%    \begin{macrocode}
\tl_new:N \g_@@_row_tl
%    \end{macrocode}
%  \end{macro}


%  \begin{macro}{\@@_handle_slot:n,\@@_handle_slot:x}
%    If the current slot number under inspection contains a glyph in
%    our font we want to typeset it. But we don't do this immediately,
%    instead we build up the whole row and typeset it later.  We
%    therefore append a \verb=&= and the glyph (including the necessary
%    formatting) to the token list \cs{g_@@_row_tl}.
%    \begin{macrocode}
\cs_new:Npn \@@_handle_slot:n #1 {
  \@@_if_uchar_exists:nTF { #1 } {
    \tl_gput_right:Nn \g_@@_row_tl
    { & \@@_format_glyph:n { \symbol{#1} } }
%    \end{macrocode}
%    We then increment the overall glyph count and record that we have
%    seen at least one glyph in the current row. There is not much
%    point in displaying rows that are completely empty; indeed,
%    we'd end up with extremely large tables which are
%    mostly empty.
%    \begin{macrocode}
    \int_gincr:N\g_@@_glyph_int
    \bool_gset_true:N \g_@@_glyph_seen_bool
%    \end{macrocode}
%    If we do font comparison we also check if the glyph is in the
%    second font and if so record that fact.
%    \begin{macrocode}
    \tl_if_empty:NF \l_@@_compare_font_tl {
      \group_begin:
        \l_@@_compare_font_tl
        \@@_if_uchar_exists:nT { #1 }
        { \int_gincr:N \g_@@_glyph_also_B_int }
      \group_end:
    }
  }
%    \end{macrocode}
%    If the current slot has no glyph in the font we also add a
%    \verb=&= followed by something that indicates the glyph is
%    missing.  If we do font comparison, it may show the glyph from
%    the other font (if it exists there) in some special way to indicate which glyph
%    should be in this slot.
%    \begin{macrocode}
  { \@@_handle_missing_glyph:n {#1} }
}
%    \end{macrocode}
%    
%    \begin{macrocode}
\cs_generate_variant:Nn \@@_handle_slot:n {x}
%    \end{macrocode}
%  \end{macro}
%

%  \begin{macro}{\@@_handle_missing_glyph:n,
%                \@@_handle_missing_glyph_std:n,
%                \@@_handle_missing_glyph_compare:n}
%    
%    In the standard case we typeset a special symbol to indicate that the glyph is missing.
%    For this case we provide some customization through keys:
%    \cs{l_@@_missing_glyph_tl} holds the symbol for a missing glyph
%    (default: a hyphen). It is typeset in a specific color and we allow for
%    setting it in a special font. The actual symbol number in
%    \verb=#1= is not needed in this scenario.
%    \begin{macrocode}
\cs_new:Npn \@@_handle_missing_glyph_std:n #1 {
  \tl_gput_right:Nn \g_@@_row_tl 
  { &
    \@@_format_glyph:n {
%        \colorbox{black!30}                      % <--- povide interface
      {\l_@@_missing_glyph_color_tl
        \l_@@_missing_glyph_font_tl 
      \l_@@_missing_glyph_tl }
    }
  }
}
%    \end{macrocode}
%
% \keysetup{missing glyphs}
%    Here are the keys for customizing the missing glyph representation.
%    \begin{macrocode}
\keys_define:nn {@@} {
   missing-glyph-color .choice:
  ,missing-glyph-color / none    .code:n    =
       \tl_clear:N \l_@@_missing_glyph_color_tl
  ,missing-glyph-color / unknown .code:n    =
       \tl_set:Nn \l_@@_missing_glyph_color_tl { \color {#1} }
%
  ,missing-glyph-font  .tl_set:N  = \l_@@_missing_glyph_font_tl
  ,missing-glyph-font  .initial:n = \ttfamily \scriptsize
  ,missing-glyph       .tl_set:N  = \l_@@_missing_glyph_tl
  ,missing-glyph       .initial:n = -  }
%    \end{macrocode}
%
%
%    The default definition for the color is to use the same as the
%    one specified by the \key{color} key. We therefore define the
%    default outside of the \pkg{l3keys} method.
%    \begin{macrocode}
\tl_new:N  \l_@@_missing_glyph_color_tl
\tl_set:Nn \l_@@_missing_glyph_color_tl {\l_@@_color_tl}
%    \end{macrocode}
%
%    This is the version that handles a missing glyph by checking
%    the \key{compare-with} font to see if that font contains the
%    glyph.
%    If yes, the substitute glyph will be typeset, otherwise the missing
%    glyph symbol is shown by calling \cs{@@_handle_missing_glyph_std:n}.
%    \begin{macrocode}
\cs_new:Npn \@@_handle_missing_glyph_compare:n #1 {
  \group_begin:
%    \end{macrocode}
%    Locally switch  to the other font, then check for the glyph:
%    \begin{macrocode}
    \l_@@_compare_font_tl
    \@@_if_uchar_exists:nTF { #1 } {
%    \end{macrocode}
%    If available, format it (together with the \texttt{\&}) but use a
%    special color and perhaps a background color.
%    \begin{macrocode}
      \tl_gput_right:Nn \g_@@_row_tl 
      { &
        \@@_format_glyph:n
        { \l_@@_compare_bgcolor_tl { \l_@@_compare_color_tl
            \l_@@_compare_font_tl
          \symbol {#1} } 
        }
      }
%    \end{macrocode}
%    Having seen a glyph only in the second font we record this fact.
%    \begin{macrocode}
      \int_gincr:N \g_@@_glyph_only_B_int            
%    \end{macrocode}
%    Also tell the algorithm that we have seen a glyph to typeset. If
%    we don't do this then a row consisting of only substitute glyphs is not
%    typeset. However, we don't update the glyph count, because this
%    is not a glyph from the main font we display.
%    \begin{macrocode}
      \bool_gset_true:N \g_@@_glyph_seen_bool
    }
%    \end{macrocode}
%    If the alternate font doesn't have the glyph either we
%    typeset the missing glyph symbol.
%    \begin{macrocode}
    { \@@_handle_missing_glyph_std:n {} }
  \group_end:
}
%    \end{macrocode}
%    
% \keysetup{comparison}
%
%    In order to display glyphs from a secondary font we need a
%    secondary color for the glyph itself and possibly some background color.
%    \begin{macrocode}
\tl_new:N \l_@@_compare_with_tl
\tl_new:N \l_@@_compare_color_tl
\tl_new:N \l_@@_compare_bgcolor_tl
%    \end{macrocode}
%    
%    \begin{macrocode}
\keys_define:nn {@@}
{
  ,compare-with .tl_set:N  = \l_@@_compare_with_tl
  ,compare-with .initial:n = 
  ,compare-color .choice:
  ,compare-color / none    .code:n
  = \tl_clear:N \l_@@_compare_color_tl
  ,compare-color / unknown .code:n
  = \tl_set:Nn  \l_@@_compare_color_tl { \color {#1} }
  ,compare-color           .initial:n = red
  ,compare-bgcolor .choice:
  ,compare-bgcolor / none    .code:n
  = \tl_clear:N \l_@@_compare_bgcolor_tl
  ,compare-bgcolor / unknown .code:n
  = \tl_set:Nn  \l_@@_compare_bgcolor_tl { \colorbox {#1} }
  ,compare-bgcolor           .initial:n = black!10
%    \end{macrocode}
%    If we run a comparison we show different statistics that have
%    their own key.
% \changes{v1.0g}{2022/11/12}{Change default text so that it makes
%    more sense if only a portion of the font is displayed (gh/4)}
%    \begin{macrocode}
  ,statistics-compare-format  .cs_set:Np
  = \@@_format_compare_stats:nnnnnn #1#2#3#4#5#6
  ,statistics-compare-format  .initial:n
  = \parbox{\linewidth}{
    Total~ number~ of~ glyphs~ shown~ from~ \texttt{#1}:~#2\\
    Comparison~ font~ \texttt{#3}~ has~ #5~ missing~ and~ #6~
  extra~ glyphs}
}
%    \end{macrocode}
%
%
%    By default, i.e., if no font for comparison has been specified, we
%    handle missing glyphs by displaying a missing glyph symbol.
%    \begin{macrocode}
\cs_new_eq:NN \@@_handle_missing_glyph:n \@@_handle_missing_glyph_std:n
%    \end{macrocode}
%  \end{macro}
%


%  \begin{macro}{\@@_format_glyph:n}
%    Every glyph is typeset in a box of equal width with the glyph
%    centered and if necessary protruding on both sides.
%    \begin{macrocode}
\cs_new:Npn \@@_format_glyph:n #1 {
  \hbox_to_wd:nn {\l_@@_glyph_box_dim} { \hss #1 \hss } }
%    \end{macrocode}
%    
% \keysetup{glyph typesetting}
%    The key to customize the width. The 6pt are fine for most cases.
%    \begin{macrocode}
\dim_new:N\l_@@_glyph_box_dim
%    \end{macrocode}
%
%    \begin{macrocode}
\keys_define:nn {@@} {
   glyph-width .dim_set:N = \l_@@_glyph_box_dim
  ,glyph-width .initial:n = 6pt
}
%    \end{macrocode}
%  \end{macro}


%  \begin{macro}{\@@_if_uchar_exists:n}
%    For testing whether or not a slot position contains a glyph we
%    need to resort to low-level methods, because so far
%    \texttt{expl3} doesn't offer an interface.
%    \begin{macrocode}
\prg_set_conditional:Npnn \@@_if_uchar_exists:n #1 { TF , T }
  { \tex_iffontchar:D \tex_font:D #1 \scan_stop:
      \prg_return_true:
    \else:
      \prg_return_false:
    \fi:
  }
%    \end{macrocode}
%  \end{macro}
%
%
%
%
% \subsection{Handling a single row}
%
%
%  \begin{macro}{\@@_handle_hex_C:n}
%    As promised here is the read definition for
%    \cs{@@_handle_hex_C:n} in all its glory.
%    \begin{macrocode}
\cs_set:Npn \@@_handle_hex_C:n #1 {
%    \end{macrocode}
%    We are now at the start of a new row (but with the last row not
%    yet typeset) and this last row may need a Unicode block heading
%    before it. This is the reason why we have to delay the
%    typesetting, because in case the line doesn't contain any glyphs
%    we want to typeset neither and that is only known after all
%    slots in the row have been processed.
%    \begin{macrocode}
  \@@_maybe_typeset_a_row_and_display_a_block_title:
%    \end{macrocode}
%    We then store away the value for the third hex digit (denoted as
%    C) in order to start with the next row.
%    \begin{macrocode}
  \tl_gset:Nn\g_@@_hex_C_tl{#1}
%    \end{macrocode}
%    Being at the start of a new row we might be at the start of a new
%    Unicode block. If so we have to update the block title to add in
%    front of the row when we typeset it (or in front of one of the
%    next rows if the first rows in the is block have no glyphs). If
%    we are still in the same block no update happens.
%    \begin{macrocode}
  \@@_update_block_title:n { \g_@@_hex_H_tl
                             \g_@@_hex_A_tl
                             \g_@@_hex_B_tl
                             \g_@@_hex_C_tl }
%    \end{macrocode}
%    We now check if this row is within the requested range, i.e.,
%    greater than or equal to \cs{l_@@_range_start_tl} and not greater than
%    \cs{l_@@_range_end_tl}.
%    \begin{macrocode}
  \int_compare:nNnF { 
    " \g_@@_hex_H_tl \g_@@_hex_A_tl
    \g_@@_hex_B_tl \g_@@_hex_C_tl 0 
  }
  < { "\l_@@_range_start_tl }
  {
    \int_compare:nNnTF { 
      " \g_@@_hex_H_tl \g_@@_hex_A_tl
      \g_@@_hex_B_tl \g_@@_hex_C_tl 0 
    }
    > { "\l_@@_range_end_tl }
%    \end{macrocode}
%    If we are past the \texttt{end-range} we break out the clist
%    mapping, to avoind unnecessary repetition. This should be
%    propagated back to the outer clists as well (not done).
%    \begin{macrocode}
    {  \clist_map_break: }
%    \end{macrocode}
%    If we are within range we process the slots in the row by first
%    initializing \cs{g_@@_row_tl} with the row title (the info on the
%    left) and then loop through all slots the row to append glyphs
%    (or missing glyphs) to \cs{g_@@_row_tl} to build up everything we
%    need to finally typeset it.
%    \begin{macrocode}
    {
      \tl_gset:Nx \g_@@_row_tl
      {
        \exp_not:N \@@_format_row_title:n
        { \g_@@_hex_H_tl \g_@@_hex_A_tl
        \g_@@_hex_B_tl \g_@@_hex_C_tl  }
      }
      \clist_map_function:NN \c_@@_hex_digits_clist
      \@@_handle_hex_D:n
    }
  }
}
%    \end{macrocode}
%  \end{macro}
%
%
%  \begin{macro}{\@@_format_row_title:n}
%    The function to format the row title on the left, as used above.
% \changes{v1.0g}{2022/11/12}{Add key hex-digits-row-format to allow
%    customizing the row title on the left (gh/3)}
%    \begin{macrocode}
\cs_new:Npn \@@_format_row_title:n #1 {
  \texttt { \footnotesize \l_@@_color_tl \@@_format_row_hex_digits:n {#1} }
}
%    \end{macrocode}
%  \end{macro}
%
%
%
% \keysetup{ranges}
%
%    For the range we have two keys, its start and the end. By default the
%    whole range from 0 to FFFF is processed.
%
%    \begin{macrocode}
\tl_new:N \l_@@_range_start_tl
\tl_new:N \l_@@_range_end_tl
%    \end{macrocode}
%
%    \begin{macrocode}
\keys_define:nn {@@}
{
  ,range-start .tl_set:N  = \l_@@_range_start_tl
  ,range-start .initial:n = 0000
  ,range-end   .tl_set:N  = \l_@@_range_end_tl
  ,range-end   .initial:n = FFFF
}
%    \end{macrocode}
%
%
%
%
%
%  \begin{macro}{\@@_maybe_typeset_a_row_and_display_a_block_title:}
%    The function handles the just-finished row and, if the row does not
%    consist only of missing glyphs, typesets it. If necessary it also typesets
%    a Unicode block name first.
%    \begin{macrocode}
\cs_new:Npn \@@_maybe_typeset_a_row_and_display_a_block_title: 
{
%    \end{macrocode}
%    We first check if the row had any real glyphs.
%    \begin{macrocode}
  \bool_if:NTF \g_@@_glyph_seen_bool
    {
%    \end{macrocode}
%    If the row needs typesetting the fun part starts. We first look
%    at the content of \cs{g_@@_block_title_tl}.
%    \begin{macrocode}
      \tl_if_empty:NTF \g_@@_block_title_tl {
%    \end{macrocode}
%    It is empty we are in the middle of a block and we can ignore
%    the Unicode title. However, we have to see if the previous row
%    (or several) was missing (i.e., contained no glyphs). In that
%    case we leave a little extra space, otherwise we just finish the
%    previous row
%    \begin{macrocode}
        \bool_if:NTF \g_@@_row_missing_bool
        { \@@_debug_nl:n{A}\\[6pt] }
        { \@@_debug_nl:n{B}\\      }
      } {
%    \end{macrocode}
%    Otherwise we first have to typeset the Unicode block title (or
%    whatever should happen instead).
%    \begin{macrocode}
        \typeout{ Processing~ \tl_use:N \g_@@_block_title_tl }
        \bool_if:NTF \l_@@_display_block_bool {
%    \end{macrocode}
%    If we are to typeset the title the action depends a bit on
%    whether we are at the very first row or typesetting a later block.
%    \begin{macrocode}
          \bool_if:NTF \g_@@_first_row_bool {
            \bool_gset_false:N \g_@@_first_row_bool
            \@@_debug_nl:n{C}\\[-4pt]
          } {
            \@@_debug_nl:n{D}\\[8pt]
            \noalign{\vskip 1pt plus 1pt} % space above block: customizable?
          }
%                \noalign{\smallskip} % space above block: customizable?
          \multicolumn{17}{c}{\normalfont \bfseries
          \tl_use:N \g_@@_block_title_tl}
%    \end{macrocode}
%    After the block title is typeset we may want to add a row of hex
%    digits as well if that was requested, otherwise we only leave a
%    bit of extra space.
%    \begin{macrocode}
          \bool_if:NTF \l_@@_blockwise_hex_digits_bool
          { \@@_debug_nl:n{E}\\*
            \@@_display_row_of_hex_digits:
            \@@_debug_nl:n{H}\\*[2pt]
          } { \@@_debug_nl:n{F}\\*[2pt] }
        } {
%    \end{macrocode}
%    If the Unicode block title is not typeset we may still have to do
%    someting special and again it differs if we at the very beginning
%    of the table (because there we do nothing except changing the
%    state of \cs{g_@@_first_row_bool}).
%    \begin{macrocode}
          \bool_if:NTF \g_@@_first_row_bool
          { \bool_gset_false:N \g_@@_first_row_bool }
          {
            \@@_debug_nl:n{G~ (new~ block)}
            \l_@@_display_block_action_tl
          }
        }
%    \end{macrocode}
%    Once we are past the block title we clear it, so that it is not
%    retypeset before the next row.
%    \begin{macrocode}
        \tl_gclear:N \g_@@_block_title_tl
      }
%    \end{macrocode}
%    The final action is to typeset the row and reset the booleans (in
%    case they were true; if they are false already then we do this
%    unnecessarily, but that is probably faster than testing first).
%    \begin{macrocode}
      \bool_gset_false:N \g_@@_glyph_seen_bool     
      \bool_gset_false:N \g_@@_row_missing_bool
      \tl_use:N \g_@@_row_tl
    }
%    \end{macrocode}
%    Current row had no glyphs; remember that fact, and that is all we
%    have to do in that case.
%    \begin{macrocode}
    {
      \bool_gset_true:N \g_@@_row_missing_bool
    }
}
%    \end{macrocode}
%  \end{macro}

%
%
% \subsection{Initialisation at the start of the table}
%
%
%  \begin{macro}{\g_@@_first_row_bool,\g_@@_glyph_seen_bool,\g_@@_row_missing_bool}
%    Declare the three booleans used in the code below. They will tell
%    us answers to the following questions:
%    \begin{itemize}
%    \item Are we processing the first row?
%    \item Have we seen any glyph so far (in the current row)?
%    \item Did we have one or more missing rows recently?
%    \end{itemize}
%    \begin{macrocode}
\bool_new:N \g_@@_first_row_bool
\bool_new:N \g_@@_glyph_seen_bool
\bool_new:N \g_@@_row_missing_bool
%    \end{macrocode}
%  \end{macro}






%  \begin{macro}{\@@_initialize_table_rows:}
%    At the start of a table we are processing the first row
%    and so we (obviously) haven't seen a glyph yet and there wasn't a
%    missing row recently.
%    \begin{macrocode}
\cs_new:Npn \@@_initialize_table_rows: {
  \bool_gset_true:N \g_@@_first_row_bool
  \bool_gset_false:N \g_@@_glyph_seen_bool     
  \bool_gset_false:N \g_@@_row_missing_bool
%    \end{macrocode}
%    And clearly the glyph count for the font(s) is zero.
%    \begin{macrocode}
  \int_gzero:N \g_@@_glyph_int 
  \int_gzero:N \g_@@_glyph_only_B_int 
  \int_gzero:N \g_@@_glyph_also_B_int 
}
%    \end{macrocode}
%  \end{macro}
%
%
%
%
% \subsection{Handling block titles}
%
%  \begin{macro}{g_@@_block_title_tl}
%    We keep the current block title in this token list.
%    \begin{macrocode}
\tl_new:N \g_@@_block_title_tl
%    \end{macrocode}
%  \end{macro}

  
%  \begin{macro}{\@@_update_block_title:n}
%    A block title is updated when the hex digits A,B,C have a certain
%    value, so this is nothing more than a huge case switch.
%    \begin{macrocode}
\cs_new:Npn \@@_update_block_title:n #1 {
  \tl_gset:Nx \g_@@_block_title_tl  {
    \int_case:nnF{ "#1 } {
      { "000 }{ Basic~ Latin }
      { "008 }{ Latin-1~ Supplement }
      { "010 }{ Latin~ Extended-A }
      { "018 }{ Latin~ Extended-B }
      { "025 }{ IPA~ Extensions }
      { "02B }{ Spacing~ Modifier~ Letters }
      { "030 }{ Combining~ Diacritical~ Marks }
      { "037 }{ Greek~ and~ Coptic }
      { "040 }{ Cyrillic }
      { "053 }{ Armenian }
      { "059 }{ Hebrew }
      { "060 }{ Arabic }
      { "070 }{ Syriac }
      { "075 }{ Arabic~ Supplement }
      { "078 }{ Thaana }
      { "07C }{ NKo }
      { "090 }{ Devanagari }
      { "098 }{ Bengali }
      { "0A0 }{ Gurmukhi }
      { "0A8 }{ Gujarati }
      { "0B0 }{ Oriya }
      { "0B8 }{ Tamil }
      { "0C0 }{ Telugu }
      { "0C8 }{ Kannada }
      { "0D0 }{ Malayalam }
      { "0D8 }{ Sinhala }
      { "0E0 }{ Thai }
      { "0E8 }{ Lao }
      { "0F0 }{ Tibetan }
      { "100 }{ Myanmar }
      { "10A }{ Georgian }
      { "110 }{ Hangul~ Jamo }
      { "120 }{ Ethiopic  }
      { "138 }{ Ethiopic~ Supplement }
      { "13A }{ Cherokee }
      { "140 }{ Unified~ Canadian~ Aboriginal~ Syllabics }
      { "168 }{ Ogham }
      { "16A }{ Runic }
      { "170 }{ Tagalog }
      { "172 }{ Hanunoo }
      { "174 }{ Buhid }
      { "176 }{ Tagbanwa }
      { "178 }{ Khmer }
      { "180 }{ Mongolian }
      { "190 }{ Limbu }
      { "195 }{ Tai~ Le }
      { "198 }{ New~ Tai~ Le }
      { "19E }{ Khmer~ Symbols }
      { "1A0 }{ Buginese }
      { "1B0 }{ Balinese }
      { "1D0 }{ Phonetic~ Extensions }
      { "1D8 }{ Phonetic~ Extensions~ Supplement }
      { "1DC }{ Combining~ Diacritical~ Marks~ Supplement }
      { "1E0 }{ Latin~ Extended~ Additional }
      { "1F0 }{ Greek~ Extended }
      { "200 }{ General~ Punctuation }
      { "207 }{ Superscripts~ and~ Subscripts }
      { "20A }{ Currency~ Symbols }
      { "20D }{ Combining~ Diacritical~ Marks~ for~ Symbols }
      { "210 }{ Letterlike~ Symbols }
      { "215 }{ Number~ Forms }
      { "219 }{ Arrows }
      { "220 }{ Mathematical~ Operators }
      { "230 }{ Miscellaneous~ Technical }
      { "240 }{ Control~ Pictures }
      { "244 }{ Optical~ Character~ Recognition }
      { "246 }{ Enclosed~ Alphanumerics }
      { "250 }{ Box~ Drawing }
      { "258 }{ Block~ Elements }
      { "25A }{ Geometric~ Shapes }
      { "260 }{ Miscellaneous~ Shapes }
      { "270 }{ Dingbats }
      { "27C }{ Miscellaneous~ Mathematical~ Symbols-A }
      { "27F }{ Supplemental~ Arrows-A }
      { "280 }{ Braille~ Patterns }
      { "290 }{ Supplemental~ Arrows-B }
      { "298 }{ Miscellaneous~ Mathematical~ Symbols-B }
      { "2A0 }{ Supplemental~ Mathematical~ Operators }
      { "2B0 }{ Miscellaneous~ Symbols~ and~ Arrows }
      { "2C0 }{ Glagolitic }
      { "2C6 }{ Latin~ Extended-C }
      { "2C8 }{ Coptic }
      { "2D0 }{ Georgian~ Supplement }
      { "2D3 }{ Tifinagh }
      { "2D8 }{ Ethiopic~ Extended }
      { "2E0 }{ Supplemental~ Punctuation }
      { "2E8 }{ CJK~ Radicals~ Supplement }
      { "2F0 }{ Kangxi~ Radicals }
      { "2FF }{ Ideographic~ Description~ Characters }
      { "300 }{ CJK~ Symbols~ and~ Punctuation }
      { "304 }{ Hiragana }
      { "30A }{ Katakana }
      { "310 }{ Bopomofo }
      { "313 }{ Hangul~ Compatibility~ Jamo }
      { "319 }{ Kanbun }
      { "31A }{ Bopomofo~ Extended }
      { "31C }{ CJK~ Strokes }
      { "31F }{ Katakana~ Phonetic~ Extensions }
      { "320 }{ Enclosed~ CJK~ Letters~ and~ Months }
      { "330 }{ CJK~ Compatibility }
      { "4DC }{ Yijing~ Hexagram~ Symbols }
      { "A00 }{ Yi~ Syllables }
      { "A49 }{ Yi~ Radicals }
      { "A70 }{ Modifier~ Tone~ Letters }
      { "A72 }{ Latin~ Extended-D }
      { "A80 }{ Syloti~ Nagri }
      { "A84 }{ Phags-pa }
      { "A88 }{ Saurashtra }
      { "A8E }{ Devanagari Extended }
      { "A90 }{ Kayah Li }
      { "A93 }{ Rejang }
      { "A96 }{ Hangul Jamo Extended-A }
      { "A98 }{ Javanese }
      { "A9E }{ Myanmar Extended-B }
      { "AA0 }{ Cham }
      { "AA6 }{ Myanmar Extended-A }
      { "AA8 }{ Tai Viet }
      { "AAE }{ Meetei Mayek Extensions }
      { "AB0 }{ Ethiopic Extended-A }
      { "AB3 }{ Latin Extended-E }
      { "AB7 }{ Cherokee Supplement }
      { "ABC }{ Meetei Mayek }
      { "AC0 }{ Hangul Syllables }
      { "D7B }{ Hangul Jamo Extended-B }
      { "D80 }{ High Surrogates }
      { "DB8 }{ High Private Use Surrogates }
      { "DC0 }{ Low Surrogates }
      { "E00 }{ Private~ Use~ Area }
      { "F90 }{ CJK~ Compatibility~ Ideographs }
      { "FB0 }{ Alphabetic~ Presentation~ Forms }
      { "FB5 }{ Arabic~ Presentation~ Forms-A }
      { "FE0 }{ Variation~ Selectors }
      { "FE1 }{ Vertical~ Forms }
      { "FE2 }{ Combining~ Half~ Marks  }
      { "FE3 }{ CJK~ Compatibility~ Forms }
      { "FE5 }{ Small~ Form~ Variants }
      { "FE7 }{ Arabic~ Presentation~ Forms-B }
      { "FF0 }{ Halfwidth~ and Fullwidth~ Forms }
      { "FFF }{ Specials~ ... }
%% ... Plane 1 ...         
      { "1000 }{ Linear~ B~ Syllabary }
      { "1008 }{ Linear~ B~ Ideograms }
      { "1010 }{ Aegean~ Numbers }
      { "1014 }{ Ancient~ Greek~ Numbers }
      { "1019 }{ Ancient~ Symbols }
      { "101D }{ Phaistos~ Disc }
      { "1028 }{ Lycian }
      { "102A }{ Carian }
      { "102E }{ Coptic~ Epact~ Numbers }
      { "1030 }{ Old~ Italic }
      { "1033 }{ Gothic }
      { "1035 }{ Old~ Permic }
      { "1038 }{ Ugaritic }
      { "103A }{ Old~ Persian }
      { "1040 }{ Deseret }
      { "1045 }{ Shavian }
      { "1048 }{ Osmanya }
      { "104B }{ Osage }
      { "1050 }{ Elbasan }
      { "1053 }{ Caucasian~ Albanian }
      { "1060 }{ Linear~ A }
      { "1080 }{ Cypriot~ Syllabary }
      { "1084 }{ Imperial~ Aramaic }
      { "1086 }{ Palmyrene }
      { "1088 }{ Nabataean }
      { "108E }{ Hatran }
      { "1090 }{ Phoenician }
      { "1092 }{ Lydian }
      { "1098 }{ Meroitic~ Hieroglyphs }
      { "109A }{ Meroitic~ Cursive }
      { "10A0 }{ Kharoshthi }
      { "10A6 }{ Old~ South~ Arabian }
      { "10A8 }{ Old~ North~ Arabian }
      { "10AC }{ Manichaean }
      { "10B0 }{ Avestan }
      { "10B4 }{ Inscriptional~ Parthian }
      { "10B6 }{ Inscriptional~ Pahlavi }
      { "10B8 }{ Psalter~ Pahlavi }
      { "10C0 }{ Old~ Turkic }
      { "10C8 }{ Old~ Hungarian }
      { "10E6 }{ Rumi~ Numeral~ Symbols }
      { "1100 }{ Brahmi }
      { "1108 }{ Kaithi }
      { "110D }{ Sora~ Sompeng }
      { "1110 }{ Chakma }
      { "1115 }{ Mahajani }
      { "1118 }{ Sharada }
      { "111E }{ Sinhala~ Archaic~ Numbers }
      { "1120 }{ Khojki }
      { "1128 }{ Multani }
      { "112B }{ Khudawadi }
      { "1130 }{ Grantha }
      { "1140 }{ Newa }
      { "1148 }{ Tirhuta }
      { "1158 }{ Siddham }
      { "1160 }{ Modi }
      { "1166 }{ Mongolian~ Supplement }
      { "1168 }{ Takri }
      { "1170 }{ Ahom }
      { "118A }{ Warang~ Citi }
      { "11A0 }{ Zanabazar~ Square }
      { "11A5 }{ Soyombo }
      { "11AC }{ Pau~ Cin~ Hau }
      { "11C0 }{ Bhaiksuki }
      { "11C7 }{ Marchen }
      { "11D0 }{ Masaram~ Gondi }
      { "1200 }{ Cuneiform }
      { "1240 }{ Cuneiform~ Numbers~ and~ Punctuation }
      { "1248 }{ Early~ Dynastic~ Cuneiform }
      { "1300 }{ Egyptian~ Hieroglyphs }
      { "1440 }{ Anatolian~ Hieroglyphs }
      { "1680 }{ Bamum~ Supplement }
      { "16A4 }{ Mro }
      { "16AD }{ Bassa~ Vah }
      { "16B0 }{ Pahawh~ Hmong }
      { "16F0 }{ Miao }
      { "16FE }{ Ideographic~ Symbols~ and~ Punctuation }
      { "1700 }{ Tangut }
      { "1880 }{ Tangut~ Components }
      { "1B00 }{ Kana~ Supplement }
      { "1B10 }{ Kana~ Extended-A }
      { "1B17 }{ Nushu }
      { "1BC0 }{ Duployan }
      { "1BCA }{ Shorthand~ Format~ Controls }
      { "1D00 }{ Byzantine~ Musical~ Symbols }
      { "1D10 }{ Musical~ Symbols }
      { "1D20 }{ Ancient~ Greek~ Musical~ Notation }
      { "1D30 }{ Tai~ Xuan~ Jing~ Symbols }
      { "1D36 }{ Counting~ Rod~ Numerals }
      { "1D40 }{ Mathematical~ Alphanumeric~ Symbols }
      { "1D80 }{ Sutton~ SignWriting }
      { "1E00 }{ Glagolitic~ Supplement }
      { "1E80 }{ Mende~ Kikakui }
      { "1E90 }{ Adlam }
      { "1EE0 }{ Arabic~ Mathematical~ Alphabetic~ Symbols }
      { "1F00 }{ Mahjong~ Tiles }
      { "1F03 }{ Domino~ Tiles }
      { "1F0A }{ Playing~ Cards }
      { "1F10 }{ Enclosed~ Alphanumeric~ Supplement }
      { "1F20 }{ Enclosed~ Ideographic~ Supplement }
      { "1F30 }{ Miscellaneous~ Symbols~ and~ Pictographs }
      { "1F60 }{ Emoticons }
      { "1F65 }{ Ornamental~ Dingbats }
      { "1F68 }{ Transport~ and~ Map~ Symbols }
      { "1F70 }{ Alchemical~ Symbols }
      { "1F78 }{ Geometric~ Shapes~ Extended }
      { "1F80 }{ Supplemental~ Arrows-C }
      { "1F90 }{ Supplemental~ Symbols~ and~ Pictographs }
      { "2000 }{ CJK~ Unified~ Ideographs~ Extension~ B }
      { "2A70 }{ CJK~ Unified~ Ideographs~ Extension~ C }
      { "2B74 }{ CJK~ Unified~ Ideographs~ Extension~ D }
      { "2B82 }{ CJK~ Unified~ Ideographs~ Extension~ E }
      { "2CEB }{ CJK~ Unified~ Ideographs~ Extension~ F }
      { "2F80 }{ CJK~ Compatibility~ Ideographs~ Supplement }
      { "E010 }{ Tags }
      { "E000 }{ Variation~ Selectors~ Supplement }
      { "F000 }{ Supplementary~ Private~ Use~ Area-A }
% higher up not covered!
    }
%    \end{macrocode}
%    If none of the above has matched we are somewhere within a block
%    so we want keep the current name. However, since the case
%    statement was executed within a |\tl_gset:Nx| we have to do this
%    by passing the current block name back.
%    \begin{macrocode}
    { \tl_use:N \g_@@_block_title_tl }
  }
}
%    \end{macrocode}
%  \end{macro}
%    
%
% \keysetup{display blocks}
%
%    The Unicode blocks may get indicated in different ways: with titles,
%    only through rules, or not at all. Here is the necessary setup.
%
%    \begin{macrocode}
\bool_new:N \l_@@_display_block_bool
\tl_new:N   \l_@@_display_block_action_tl
%    \end{macrocode}
%
%    \begin{macrocode}
\keys_define:nn {@@}
{
  ,display-block .choice:
  ,display-block / titles  .code:n =
  \bool_set_true:N \l_@@_display_block_bool
  \tl_set:Nn \l_@@_display_block_action_tl {\\}
  ,display-block / rules   .code:n =
  \bool_set_false:N \l_@@_display_block_bool
  \tl_set:Nn \l_@@_display_block_action_tl {\\ \midrule}
  ,display-block / none   .code:n =
  \bool_set_false:N \l_@@_display_block_bool
  \tl_set:Nn \l_@@_display_block_action_tl {\\}
  ,display-block .initial:n  = titles
}
%    \end{macrocode}
%
%    That's all of the programming using the L3 layer.
%    \begin{macrocode}
\ExplSyntaxOff
%    \end{macrocode}
%    What remains is to require all packages needed \ldots
%    \begin{macrocode}
\RequirePackage{longtable,booktabs,caption}
%    \end{macrocode}
%    \ldots and executing all options passed to the
%    package via \cs{usepackage}.
%    \begin{macrocode}
\ProcessKeysPackageOptions{@@}
%</package>
%    \end{macrocode}
%
% \Finale
\endinput

